%% 
%% Copyright 2007-2025 Elsevier Ltd
%% 
%% This file is part of the 'Elsarticle Bundle'.
%% ---------------------------------------------
%% 
%% It may be distributed under the conditions of the LaTeX Project Public
%% License, either version 1.3 of this license or (at your option) any
%% later version.  The latest version of this license is in
%%    http://www.latex-project.org/lppl.txt
%% and version 1.3 or later is part of all distributions of LaTeX
%% version 1999/12/01 or later.
%% 
%% The list of all files belonging to the 'Elsarticle Bundle' is
%% given in the file `manifest.txt'.
%% 
%% Template article for Elsevier's document class `elsarticle'
%% with numbered style bibliographic references
%% SP 2008/03/01
%% $Id: elsarticle-template-num.tex 272 2025-01-09 17:36:26Z rishi $
%%
\documentclass[preprint,12pt]{elsarticle}
%% Use the option review to obtain double line spacing
%% \documentclass[authoryear,preprint,review,12pt]{elsarticle}

%% Use the options 1p,twocolumn; 3p; 3p,twocolumn; 5p; or 5p,twocolumn
%% for a journal layout:
%% \documentclass[final,1p,times]{elsarticle}
%% \documentclass[final,1p,times,twocolumn]{elsarticle}
%% \documentclass[final,3p,times]{elsarticle}
%% \documentclass[final,3p,times,twocolumn]{elsarticle}
%% \documentclass[final,5p,times]{elsarticle}
%% \documentclass[final,5p,times,twocolumn]{elsarticle}

%% For including figures, graphicx.sty has been loaded in
%% elsarticle.cls. If you prefer to use the old commands
%% please give \usepackage{epsfig}
%% The amssymb package provides various useful mathematical symbols
\usepackage{amssymb}
%% The amsmath package provides various useful equation environments.
\usepackage{amsmath}
\usepackage{algorithm}
\usepackage{algorithmicx}
\usepackage{algpseudocode} % o algorithmic
\usepackage{url}
\usepackage{booktabs}
\usepackage{multirow}
\usepackage{float}
\floatstyle{plaintop}
\restylefloat{table}
\usepackage{graphicx}
\usepackage{verbatim}
\usepackage{subcaption}
\usepackage{xcolor}
%% The amsthm package provides extended theorem environments
%% \usepackage{amsthm}
\usepackage[labelfont=bf,labelsep=none]{caption} % controla estilo de captions
\renewcommand{\figurename}{\textbf{Fig.}}


% Cambiar "Table" por "Tabla"
\renewcommand{\tablename}{Table}

% Poner el caption arriba de la tabla
\captionsetup[table]{position=top,justification=raggedright,singlelinecheck=false}

% Quitar punto final automático (por si lo añade)
\captionsetup{labelsep=none}
% Cambiar "Figure" por "Fig." en negrita


%% The lineno packages adds line numbers. Start line numbering with
%% \begin{linenumbers}, end it with \end{linenumbers}. Or switch it on
%% for the whole article with \linenumbers.
%% \usepackage{lineno}

\journal{Information Sciences}

\begin{document}

\begin{frontmatter}

%% Title, authors and addresses

%% use the tnoteref command within \title for footnotes;
%% use the tnotetext command for theassociated footnote;
%% use the fnref command within \author or \affiliation for footnotes;
%% use the fntext command for theassociated footnote;
%% use the corref command within \author for corresponding author footnotes;
%% use the cortext command for theassociated footnote;
%% use the ead command for the email address,
%% and the form \ead[url] for the home page:
%% \title{Title\tnoteref{label1}}
%% \tnotetext[label1]{}
%% \author{Name\corref{cor1}\fnref{label2}}
%% \ead{email address}
%% \ead[url]{home page}
%% \fntext[label2]{}
%% \cortext[cor1]{}
%% \affiliation{organization={},
%%             addressline={},
%%             city={},
%%             postcode={},
%%             state={},
%%             country={}}
%% \fntext[label3]{}

\title{EDRS: Extremity-Density Representative Selection for Semi-Supervised Learning on Imbalanced Data}

%% Author 1
\author[1]{Alberto Durán-López}
\ead{albduranlopez@ugr.es}

%% Author 2 (corresponding)
\author[1]{Daniel Bolaños-Martinez\corref{cor1}}
\ead{danibolanos@ugr.es}

%% Author 3
\author[1,2]{Maria Bermudez-Edo}
\ead{mbe@ugr.es}

\cortext[cor1]{Corresponding author.}

%% Affiliation 1
\affiliation[1]{organization={Research Centre for Information and Communication Technologies (CITIC-UGR), University of Granada},
            city={Granada},
            postcode={18014},
            country={Spain}}

%% Affiliation 2
\affiliation[2]{organization={Andalusian Institute of Research and Innovation in Tourism (IATUR)},
            city={Granada},
            postcode={18011},
            country={Spain}}

\begin{abstract}
%% Text of abstract 200 palabras máximo
Representative sample selection improves training in semi-supervised learning (SSL) where labeled data are limited and must reflect the original dataset. Recent SSL methods ignore class imbalance and lack tabular data case studies. To fill this gap, we propose Extremity-Density Representative Selection (EDRS), a preprocessing point selection method for imbalanced tabular datasets. EDRS ranks unlabeled candidates combining two scores: density, which favors regions with many individuals, and extremity, which ensures inclusion of extreme cases likely belonging to minority classes. We first cluster the data to ensure diverse and representative coverage of the space, and then select samples with the highest density and extremity values, balancing outlier avoidance with coverage of extreme values. EDRS is used to select samples for labeling in a SSL framework and is compared with Random Sampling, Stratified Sampling, K-Means–derived methods, USL, Hybrid-CEAL, FDMat, Gaussian Mapping, and ESC-FFS. We validate EDRS on twelve synthetic and six real-world imbalanced datasets using SSL VIME, Manifold Mixup and Contrastive Mixup. EDRS achieves a class imbalance ratio (IR) close to 1 {\color{blue}and is} 99\% faster than other algorithms with similar IR, improves F1-score by 3–5\% in well-separated classes, and includes an ablation test evaluating the impact of density and extremity.
\end{abstract}

%%Research highlights
\begin{highlights}
\item EDRS merges density and extremity for representative sample selection in tabular SSL

\item Achieves near-balanced selections (IR$\sim$1) when classes are well separated (FDR$\geq$1)

\item Improves F1-score by 3–5\% while reducing selection time by 90-99\% in SSL tasks

\item Validated on 12 synthetic and 6 real tabular datasets with imbalanced binary classes
\end{highlights}

%% Keywords entre 1 y 7 palabras
\begin{keyword}
Semi-supervised learning \sep Internet of Things (IoT) \sep Imbalanced classes \sep Representative sample selection \sep Tabular data
%% keywords here, in the form: keyword \sep keyword

%% PACS codes here, in the form: \PACS code \sep code

%% MSC codes here, in the form: \MSC code \sep code
%% or \MSC[2008] code \sep code (2000 is the default)

\end{keyword}

\end{frontmatter}

%% Add \usepackage{lineno} before \begin{document} and uncomment 
%% following line to enable line numbers
%% \linenumbers

%% main text
%%

\section{Introduction}

Representative sample selection consists of choosing a subset of samples contained in a larger dataset that reflects the underlying data distribution \cite{liebman2015representative}. Researchers have applied this process to different machine learning (ML) tasks, including clustering or classification, as a way to reduce computational cost while preserving performance \cite{olvera2010review,chen2025fpsmix}. Instead of training models on the entire dataset, which can be expensive and redundant, representative selection focuses on the most informative samples \cite{olvera2010review}. Given the benefits observed, its use has extended to semi-supervised learning (SSL) \cite{wang2022unsupervised,roy2024impact}. In SSL, first, algorithms learn class boundaries from a labeled dataset. Second, {\color{blue}they} exploit patterns in unlabeled data (e.g. clusters, correlated features, or density regions) to refine the model and improve those boundaries \cite{van2020survey,cai2023review}. The selection of representative samples is important for SSL, especially when the annotation cost (or budget) is limited, since labeling requires manual effort \cite{you2020self,singh2024shifting}. If the labeled examples are not representative of the overall data distribution, the most effective algorithm may perform poorly \cite{singh2024shifting}. 

In SSL, labeled data can be formed either by random sampling, where selected samples may not be representative \cite{berthelot2019remixmatch}, or by stratified sampling, which assumes access to all true labels \cite{wang2022debiased}. Ideally, samples should be chosen from unlabeled data using a method that identifies the most representative ones without label information. Prior work explores unsupervised methods, such as clustering, to achieve this \cite{wang2022unsupervised}. For example, Roy and Etemad use clustering methods like K-Means to group data and select representative samples from cluster centroids to increase diversity \cite{roy2024impact}. Other works add a density score before clustering to prioritize diverse samples with high local density \cite{wang2022unsupervised,wang2025resampling}. However, these general unsupervised approaches often struggle with tabular structures and do not inherently address the skewed distributions typical of imbalanced datasets.


Recently, methods relevant to tabular domains have emerged, such as hybrid cost-effective active learning (Hybrid-CEAL), which integrates outlier detection to capture diverse samples \cite{shankaranarayana2024hybrid}. Furthermore, data-agnostic approaches adaptable to tabular data \cite{hadar2024datamap}, such as the core-set selection method feature distribution matching (FDMat), employ optimal transport to align the subset feature distribution with the entire dataset \cite{xiao2024feature}. To tackle the class imbalance problem, methods have been proposed in the image field. For instance, the farthest-first search (FFS) algorithm iteratively chooses exemplars that cover all dataset subspaces to ensure minority class coverage \cite{you2020self}. Similarly, representation kernel alignment (RKA) utilizes Gaussian kernels to map unlabeled data towards labeled class prototypes, effectively leveraging density information to refine minority pseudo-labels \cite{li2025kernel}. However, a gap remains: tabular-data methods rarely address class imbalance, while techniques that deal with imbalance are mainly designed for image data.

To address this gap, we propose the Extremity-Density Representative Selection (EDRS) algorithm, which identifies a representative subset from unlabeled data. EDRS evaluates two scores per sample: density, a proven criterion for selecting representative samples, and extremity, a metric that balances class representation in imbalanced binary problems. Following Wang et al.'s work, we apply clustering to select, in an unsupervised manner, the samples with the highest scores \cite{wang2022unsupervised}. In our case, we select samples with the highest combined score for density and extremity, a composite metric specifically designed to optimize the trade-off between representativeness and minority class coverage in imbalanced datasets. This combination addresses an important limitation in existing methods: while standard density-based SSL primarily focuses on representativeness (exploitation), often leading to a bias towards the majority class \cite{wang2022debiased, li2024class}, traditional active learning methods typically rely on uncertainty (exploration), which can be susceptible to selecting noisy instances or outliers \cite{li2024active, kasarla2024maximally}. EDRS overcomes this by integrating a preliminary Local Outlier Factor (LOF) filtration with a hybrid extremity-density selection. This ensures that the exploration of minority class boundaries is performed only on reliable data, preventing the noisy outlier problem common in diversity-based baselines \cite{shao2024enhancing, xiao2024feature}.

The main contributions of this paper are summarized as follows: 

\begin{itemize}
    \item The proposal of EDRS, a representative sample selection method that integrates density and extremity criteria for SSL. Specifically, our method computes a density score and an extremity score for each unlabeled instance, then chooses samples with a high combined score. 

    \item The creation of twelve synthetic imbalanced binary tabular datasets for evaluation purposes. We also include the creation of our in-house smart‑village dataset.

    \item A comprehensive comparison of the proposed method against the state-of-the-art methods. EDRS achieves better results than the evaluated state-of-the-art methods across multiple SSL algorithms.

    \item An ablation test to demonstrate the individual contribution of density and extremity scores.
    
\end{itemize}

The remainder of this paper is organized as follows: Section~\ref{sec:related} reviews related work; Section~\ref{sec:proposed_method} outlines the proposed methodology; Section~\ref{sec:experiment} presents experimental details and discusses the results; and Section~\ref{sec:conclusion} concludes.


\section{Related Work} \label{sec:related}

Selecting representative samples for labeling improves SSL~\cite{shao2024enhancing}, motivating recent work on methods for choosing labeled subsets that improve learning. Labels are usually sampled randomly or via stratification~\cite{berthelot2019remixmatch,wang2022debiased}, but random sampling risks imbalance and poor coverage, while stratification is infeasible without true labels. Active learning methods have been tested to select representative samples~\cite{margatina2021active,ducoffe2018adversarial}, yet they require full true label knowledge during training~\cite{albelwi2016analysis}. In SSL, we must select samples from unlabeled data before annotation, demanding unsupervised point-selection methods to identify the most representative ones.

Several studies have explored unsupervised representative sample selection for SSL. Wang et al. \cite{wang2022unsupervised} propose an Unsupervised Selective Labeling (USL) method that learns feature representations via contrastive learning \cite{li2025leukocyte} and selects cluster prototypes based on density peaks and inter-cluster diversity to form a representative seed set without labels. Roy et al. \cite{roy2024impact} compare clustering approaches such as K-Means, K-Means\texttt{++}, and bisecting K-Means for grouping unlabeled samples and selecting each cluster’s centroid for labeling. They also analyze how labeling timing and batch size affect convergence and accuracy. However, these clustering-based methods were developed primarily for balanced image data.

In parallel, methods relevant to tabular data have emerged. In the context of active learning, Shankaranarayana et al.~\cite{shankaranarayana2024hybrid} introduced Hybrid-CEAL, a method specifically designed for tabular classification. This method improves representative sampling by incorporating an outlier detection module that enhances the minority classes representation. Xiao et al.~\cite{xiao2024feature} proposed FDMat, which employs optimal transport to match the subset's feature distribution with the original data. Recent works have used FDMat in tabular data, demonstrating that this methodology is data-agnostic \cite{hadar2024datamap}. However, these approaches exhibit limitations when facing imbalanced datasets. FDMat aims to replicate the global distribution. Consequently, it preserves the original class imbalance within the selected subset rather than correcting it. Similarly, Hybrid-CEAL relies on outlier detection to enhance diversity. This strategy risks treating randomly selected anomalies as representative samples, thereby confounding minority-class instances with true outliers.

To specifically tackle the class imbalance problem, other methods have been proposed. You et al. \cite{you2020self} introduce an Exemplar-based Subspace Clustering (ESC) model that frames selection as minimizing an $\ell_1$-regularized reconstruction cost under a Minkowski-functional coverage constraint, using the FFS algorithm to iteratively choose samples that best span the data subspaces. By ensuring coverage of minority-class regions, their method improves selection under imbalance. Similarly, Li et al.~\cite{li2025kernel} developed RKA, which employs Gaussian kernel mapping to align unlabeled samples with class prototypes (feature centroids). To ensure representative selection, RKA maintains dynamic centroids updated via exponential moving average (EMA), filtering out outliers that fall below a specific similarity confidence threshold. Despite these advantages, both approaches exhibit limitations. FFS prioritizes diversity, which increases the risk of selecting outliers as representative samples. Conversely, RKA targets high-density prototypes. Consequently, it potentially overlooks the boundary samples required to define decision margins, which include both dense and extremity data.

\section{Methodology} \label{sec:proposed_method}

As illustrated in Fig.~\ref{fig:scheme}, our proposal EDRS consists of six steps: outlier removal, dimensionality reduction, compute score (density and extremity score), clustering into groups, and representative sample selection. The pseudocode is provided in Algorithm~\ref{algo:edrs}, and the following subsections describe each step in detail. Steps 3, 4, and 5 can be performed in parallel, as they are independent from each other. Our methodology is optimized for numerical features, as the EDRS pipeline relies on distance-based density estimation (LOF) and variance-based dimensionality reduction (PCA), which operate in a continuous numerical space.

\begin{figure}[!t]
    \centering
    \includegraphics[width=0.85\linewidth]{Figures/proposal.png}
    \caption{. EDRS method working scheme.}\label{fig:scheme}
\end{figure}

\subsection{Outlier Removal} \label{subsec:remove_outliers}

To prevent outliers from being selected as representatives, we remove a percentage of samples (\%\_outliers in Algorithm~\ref{algo:edrs}) with the highest LOF scores~\cite{breunig2000lof}. LOF measures how isolated a sample is relative to its neighbors, detecting outliers in datasets with varying densities, and being robust to heterogeneity and global density shifts~\cite{kriegel2011interpreting}.  

To calculate the LOF, a few preliminary concepts need to be introduced first. Let $X = \{x_i\}_{i=1}^M \subset \mathbb{R}^d$ be an unlabeled dataset of size $M$. For a sample $x_1$, the reachability distance to another sample $x_2$ is
\begin{equation} \label{eq:reach_dist}
    \operatorname{reach - dist}_k(x_1,x_2) = \max\{\operatorname{k - dist}(x_2),\, d(x_1,x_2)\},
\end{equation}
where $d(x_1,x_2)$ is their distance and $\operatorname{k-dist}(x_2)$ is the distance from $x_2$ to its $k$-th nearest neighbor. The local reachability density (lrd$_k$) of $x_1$ is
\begin{equation} \label{eq:lrd}
    \operatorname{lrd}_k(x_1) = \frac{|N_k(x_1)|}{\sum_{x_2 \in N_k(x_1)} \operatorname{reach - dist}_k(x_1,x_2)},
\end{equation}
with $N_k(x_1)$ denoting its $k$ nearest neighbors. The LOF score is then
\begin{equation} \label{eq:lof}
    \operatorname{LOF}_k(x_1) = \frac{1}{|N_k(x_1)|} \cdot \sum_{x_2 \in N_k(x_1)} \frac{\operatorname{lrd}_k(x_2)}{\operatorname{lrd}_k(x_1)}.
\end{equation}
Samples with $\operatorname{LOF}_k(x_1) > 1$ are considered outliers and removed, as their local density is lower than that of their neighbors~\cite{hu2025efficient}.

\subsection{Dimensionality Reduction} \label{subsec:dimen_reduc}

To capture the main variance directions and support the extremity metric, we reduce the data to two dimensions using PCA~\cite{jolliffe2011principal}. Given $X \in \mathbb{R}^{N\times D}$, PCA finds orthonormal vectors $\mathbf{v}_1,\mathbf{v}_2 \in \mathbb{R}^D$ by solving
\begin{equation}
\mathbf{v}_1 = \arg\max_{\|\mathbf{v}\|=1}\mathrm{Var}(X\mathbf{v}), 
\quad
\mathbf{v}_2 = \arg\max_{\substack{\|\mathbf{v}\|=1\\ \mathbf{v}\perp\mathbf{v}_1}}\mathrm{Var}(X\mathbf{v}),
\end{equation}
and projects $X$ as
\begin{equation}
X_{\text{pca}} = X[\mathbf{v}_1,\mathbf{v}_2].
\end{equation}
Using two components captures most variance while reducing noise and redundancy~\cite{salem2019data}. The first component $\mathbf{v}_1$, representing the direction of maximum variance, is later used to define the extremity measure~\cite{lakhina2004diagnosing}.

It is important to clarify that while tabular data often exhibits non-linear manifold structures \cite{li2025kernel}, the primary objective of this dimensionality reduction is not to model these complex relationships, but to preserve global Euclidean distances to accurately calculate the Extremity Score. Non-linear embeddings, such as t-SNE or UMAP, prioritize the preservation of local neighborhoods, which inevitably results in the distortion of global distances. This distortion risks compressing geometrically distant outliers or minority samples into dense clusters to satisfy local topology constraints, effectively masking the extremity signal our algorithm relies on. In contrast, PCA ensures that samples geometrically distant from the center in the original Euclidean space remain distinguishable as extreme values in the projection.


\subsection{Compute Density Score} \label{subsec:density_score_calc}  

We compute a density score $d_i$ for each sample $\mathbf{x}_i \in X_{\text{pca}}$ to locate dense regions. The score is defined as the inverse of the median Euclidean distance from $\mathbf{x}_i$ to its $k'$ nearest neighbors in PCA space:

\begin{equation}
\label{eq:density_score}
d_i = \bigl(\operatorname{median}_{\mathbf{x}_j \in N_{k'}(\mathbf{x}_i)} \|\mathbf{x}_i - \mathbf{x}_j\|_2 \bigr)^{-1},
\end{equation}

where $N_{k'}(\mathbf{x}_i)$ denotes the $k'$ nearest neighbors of $\mathbf{x}_i$, and $\|\cdot\|_2$ is the $L_2$ norm, or Euclidean distance. Finally, we set $d_i = 0$ for samples below a predefined density threshold ($\delta_{\text{density}}$) to filter low-density regions.

\subsection{Compute Extremity Score} \label{subsec:extremity_composite_score}

We compute an extremity score $e_i$ for each sample $\mathbf{x}_i \in X_{\text{pca}}$ based on its position along the first principal component. Let $p_{i,1}$ be the projection of $\mathbf{x}_i$ onto $\mathbf{v}_1$. We normalize these projections to $[0,1]$ using min–max normalization:
\begin{equation} \label{eq:pc1_normalized}
    p'_{i,1} = \frac{p_{i,1} - \min_{\mathbf{x}_j \in X_{\text{pca}}}(p_{j,1})}{\max_{\mathbf{x}_k \in X_{\text{pca}}}(p_{k,1}) - \min_{\mathbf{x}_j \in X_{\text{pca}}}(p_{j,1})}
\end{equation}
where min and max are computed over all samples in $X_{\text{pca}}$. The extremity score $e_i$ is then:
\begin{equation} \label{eq:extremity_score}
    e_i = 2 \cdot \left| p'_{i,1} - 0.5 \right|.
\end{equation}
This formulation assigns scores near $1$ to samples at the extremes of the $p_{i,1}$ distribution (positives and negatives), while scores near 0 indicate samples near its center.


It is important to note that this linear definition of extremity is specifically designed for binary classification tasks, where the minority class is expected to deviate from the majority class along the principal axis of variance. For multi-class environments, where classes may be distributed radially or across multiple dimensions, a One-vs-Rest (OvR) strategy could be used to apply EDRS effectively. By iteratively isolating each class against the rest, the PCA step would dynamically re-align the principal axis ($\mathbf{v}_1$) to capture the specific variance of the target class of each iteration \cite{li2022trustworthy}, ensuring the extremity score remains valid regardless of the global spatial distribution. This strategy introduces a linear multiplier proportional to the number of classes ($C$), yielding a total complexity of $\mathcal{O}(C \cdot M \log M)$, yielding a reasonable computational overhead. The dimensionality reduction in EDRS relies on PCA, which is an analytical linear transformation computed in a single deterministic step, avoiding the quadratic $\mathcal{O}(M^2)$ complexity of iterative manifold learning methods like t-SNE. Nevertheless, to generalize EDRS beyond this iterative framework, multivariate statistical depth functions such as Tukey’s Halfspace or Mahalanobis Depth \cite{nagy2025halfspace, mozharovskyi2025anomaly} could be adopted. These methods define a global center-outward ordering in $\mathbb{R}^n$, enabling the simultaneous identification of extreme samples without relying on multiple linear projections.


\subsection{Clustering into $N$ Groups} \label{subsec:clustering_k_groups}

To ensure representative selection covers the feature space, we partition $X_{\text{pca}}$ into $N$ clusters (equal to the annotation budget). In the next step, one sample will be selected from each cluster to ensure coverage of the entire feature space. Clustering is performed with K-means++~\cite{roy2024impact,wang2022unsupervised}, whose usage speeds convergence, improves cluster consistency on large datasets, and scales linearly with data size~\cite{bachem2016fast}.

\subsection{Select Representatives with Higher Score} \label{subsec:select_representative_points}

We fuse the density $d_i$ (Equation \ref{eq:density_score}) and the extremity $e_i$ (Equation \ref{eq:extremity_score}) scores into a single composite score:
\begin{equation} \label{eq:composite_score}
    s_i = d_i \cdot \bigl(1 + e_i\bigr).
\end{equation}
The score $s_i$ highlights samples with high local density and strong extremity, favoring samples well supported by neighbors yet distinct in feature space \cite{wang2022unsupervised,qu2021outlier}. The term $(1+e_i)$ acts as a dynamic weight that amplifies the density score based on extremity.

Finally, we select a target of $N$ representative samples from the $X_{\text{pca}}$ data, utilizing the $N$ clusters obtained in Section \ref{subsec:clustering_k_groups}. For each cluster, we select the sample with the highest composite score $s_i$, balancing local density and extremity. However, as we aim to represent imbalanced datasets, the extreme values (minority class) need to be represented in the selected samples, prioritizing extremity over density. Hence, if the selected sample's extremity value $e_i$ falls below a configurable threshold ($e_{\text{thresh}}$ in Algorithm~\ref{algo:edrs}), it is replaced by the cluster member with the highest $e_i$. This ensures that each cluster contributes at least one sufficiently extreme example, improving minority-class representation in imbalanced datasets~\cite{lakhina2004diagnosing}.

\begin{algorithm}[!t]
\caption{\scriptsize EDRS: Extremity Density Representative Selection}
\label{algo:edrs}
\scriptsize            % <-- sólo afecta al bloque algorithmic
\begin{algorithmic}[1]
\Procedure{EDRS}{$\mathit{data}, N, \%\_\mathrm{outliers}, k', \delta_{\mathrm{density}}, e_{\mathrm{thresh}}$}
    \State \textbf{Step 1: Remove outliers}  
    \State $\mathit{lof\_labels} \gets \text{Apply LOF to data using \%\_outliers}$  
    \State $\mathit{clean} \gets \{\mathbf x \in \mathit{data} \mid \mathit{lof\_labels}(\mathbf x)=\text{inlier}\}$  
    \State \textbf{Step 2: PCA to 2D}  
    \State $\mathit{embeddings} \gets \text{Apply PCA to reduce data to 2 dimensions}$  
    \State \textbf{Step 3: Compute density score}
    \State $\mathit{Dists} \gets \text{Compute KNN in } embeddings \text{ using } k'$  
    \State $\mathit{density} \gets 1 / \mathrm{median}(\mathit{Dists})$  
    \State For each $i$, if $\mathit{density}_i < \delta_{\mathrm{density}}$, set $\mathit{density}_i \gets 0$  
    \State \textbf{Step 4: Compute extremity score}  
    \State $\mathit{pc_1} \gets \text{Normalize the first PCA component from } embeddings$  
    \State $\mathit{extremity} \gets 2 \times |\mathit{pc_1} - 0.5|$  
    \State \textbf{Step 5: Cluster into $N$ groups}  
    \State $\mathit{clusters} \gets \mathrm{KMeans++}(embeddings, N)$   
    \State \textbf{Step 6: Select $N$ representatives with higher score}
    \State $\mathit{score} \gets \mathit{density}\,\cdot\,(1 + \mathit{extremity})$  
    \State Initialize $\mathit{selected} \gets \emptyset$
    \For{each cluster $c=1,\dots,N$}
        \State $p^* \gets \arg\max_{\mathbf{x}\in c}\,\mathit{score}(\mathbf{x})$
        \If{$\mathit{extremity}(p^*) \le e_{\mathrm{thresh}}$}
            \State $p^* \gets \arg\max_{\mathbf{x}\in c}\,\mathit{extremity}(\mathbf{x})$
        \EndIf
        \State Add $p^*$ to $\mathit{selected}$
    \EndFor
    \State \Return $\mathit{selected}$
\EndProcedure
\end{algorithmic}
\end{algorithm}


\section{Experiments and results} \label{sec:experiment}

To evaluate our EDRS algorithm, we validate it on twelve synthetic and six real‑world datasets, comparing EDRS against ten representative sample selection algorithms from the literature. Then, we use each algorithms’ selected samples to train SSL models and evaluate their performance. This is explained in detail in the next subsection.


\subsection{Evaluation pipeline} \label{sec:method}

We introduce a pipeline to evaluate representative sample selection methods in SSL. The pipeline separates the selection method from the SSL algorithm, enabling fair comparisons across approaches. Fig. \ref{fig:pipeline} illustrates the pipeline in three phases: Representative sample selection (Phase I), data labeling (Phase II), and training and evaluation (Phase III).

\begin{figure}[!t]
    \centering
    \includegraphics[width=0.85\linewidth]{Figures/pipeline.png}
    \caption{. Overview of the proposed pipeline.}
    \label{fig:pipeline}
\end{figure}

\subsection*{Phase I: Representative sample selection}
\label{subsec:Phase1}

Let $X = \{x_i\}_{i=1}^M \subset \mathbb{R}^d$ be the unlabeled dataset of size $M$. Given an annotation budget $N$, we define an operator $S$ that selects a subset $R \subset X$ of size $|R| = N$:
\begin{equation} \label{eq:selection_operator}
  R = S(X, N)
\end{equation}
Here, $S$ denotes any representative sample selection method. The goal is for $R$ to capture the structure and distribution of $X$ while minimizing redundancy and ensuring coverage of key regions in the data space.

\subsection*{Phase II: Data Labeling}
\label{subsec:Phase2}
After obtaining $R$, each $r_j \in R$ is annotated with its true label $y_j$, forming the labeled set ($L$).
\begin{equation}
  \mathcal{L} = \{(r_j, y_j)\}_{j=1}^N
\end{equation}
As shown in Phase~II of Fig.~\ref{fig:pipeline}, only these $N$ samples are labeled. The true labels are only available for $R$, while $M - N$ samples constitute the unlabeled set ($U$).
\begin{equation}
  \mathcal{U} = X \setminus R
\end{equation}

\subsection*{Phase III: Training and Model Evaluation}\label{subsec:Phase3}

In this phase, we train the SSL model and evaluate its performance on the full dataset.

\begin{enumerate}
  \item \textbf{Model training.} A semi-supervised learner $f_{\theta}$ is trained on the labeled set $\mathcal{L}$ and unlabeled set $\mathcal{U}$ using any SSL algorithm.
  
  \item \textbf{Full-dataset prediction.} The trained model produces predictions for all samples:
  \begin{equation}
    \hat{y}_i = f_{\theta}(x_i), \quad \forall\,x_i \in X
  \end{equation}

  \item \textbf{Performance assessment.} Predictions $\hat{y}_i$ are compared with true labels from the entire dataset $X$ to compute accuracy, precision, and F1-score, evaluating how the selection method $S$ and model $f_{\theta}$ perform under annotation budget $N$. For evaluation purposes, we have access to the true values for the entire dataset, and we compare the results with the ground truth.
\end{enumerate}

These steps provide a full evaluation of each representative sample selection method within the SSL pipeline.

\subsection{Datasets}

All datasets involve binary classification with strong class imbalance. This setup allows us to assess our algorithm’s ability to select representative samples under imbalance. Table \ref{tab:datasets_summary} summarizes the datasets used in this work.

\subsection*{Synthetic Datasets}

First, we create twelve datasets, where each dataset has two classes: class 0 (simulating a dense cluster) is  
\[
\mathcal{N}\bigl((0,0)^\top,\;\Sigma_{0}\bigr),\quad
\Sigma_{0}=\begin{pmatrix}0.5\,v&0\\0&0.1\,v\end{pmatrix},
\]
and class 1 (simulating a sparse cluster) is  
\[
\mathcal{N}\bigl((\delta,0)^\top,\;\Sigma_{1}\bigr),\quad
\Sigma_{1}=\begin{pmatrix}8\,v&0\\0&1\,v\end{pmatrix},
\]
where $\mathcal{N}(\mu,\Sigma)$ denotes a multivariate normal distribution with mean vector $\mu$ and covariance matrix $\Sigma$. 
Here, $v$ scales the overall variance, and $\delta$ controls the separation between the two class means. 
Varying the total sample size, the variance factor $v$, and the separation $\delta$ simulates different cluster densities. Specifically, the first six datasets contain 800 samples each, and the remaining six contain 50,000 samples. 
Within each group, the class separation $\delta$ varies between $4$ and $0.125$ Global Fisher Discriminant Ratio (FDR), while maintaining a fixed class imbalance of 90-10.

\subsection*{Real World Datasets}

Then, we use five publicly available real world datasets from the UCI database. Breast Cancer\footnote{\url{https://archive.ics.uci.edu/dataset/17/breast+cancer+wisconsin+diagnostic}}, Thyroid Disease\footnote{\url{https://archive.ics.uci.edu/dataset/102/thyroid+disease}}, Wilt\footnote{\url{https://archive.ics.uci.edu/dataset/285/wilt}}, Seismic Bumps\footnote{\url{https://archive.ics.uci.edu/dataset/266/seismic+bumps}}, Bank Marketing\footnote{\url{https://archive.ics.uci.edu/dataset/222/bank+marketing}}.

\textbf{Smart Poqueira\footnote{\url{https://zenodo.org/records/14262136}} (own dataset)}: We also use an in-house dataset \cite{duran2025route,duran2026gct}. Data were collected over 19 months (February 2022–August 2023) from four Hikvision LPR IP cameras at the entry and exit points of a region near Sierra Nevada National Park in Granada, Spain. Each row logs a timestamp, license plate, and camera ID. We matched plates to fiscal addresses via the Spanish Directorate‑General for Traffic\footnote{\url{https://www.dgt.es/inicio/}} to identify residents. From these sources, we derived a set of variables characterizing visit patterns (e.g., total entries, duration, overnight stays), temporal context based on the national calendar (e.g., holidays, high/low season), and the ground-truth residency label. The dataset comprises 80,750 individuals, of whom 80,006 ($\sim$99\%) are visitors to the area and 744 ($\sim$1\%) are residents.
    

\begin{table}[!t]
\centering
\caption{\\ Summary of the datasets, including class balance, size, and a brief description}
\label{tab:datasets_summary}
\resizebox{\textwidth}{!}{%
\begin{tabular}{llll}
\toprule
\textbf{Dataset} & \textbf{Balance Classes} & \textbf{Size} & \textbf{\parbox{5cm}{Description}} \\
\midrule

Synthetic Datasets 
& 90\% class 0 / 10\% class 1 
& 800; 50,000 
& \parbox{7cm}{Generated from two bivariate normal distributions one dense and the other sparse.\vspace{0.5em}} \\

Breast Cancer 
& 63\% benign / 37\% malignant 
& 569 
& \parbox{7cm}{Diagnostic features extracted from digitized images of breast tissue.\vspace{0.5em}} \\

Thyroid Disease 
& 92\% normal / 8\% disorder 
& 3772 
& \parbox{7cm}{Contains clinical and laboratory values such as T3, T4, and TSH hormone levels.\vspace{0.5em}} \\

Wilt 
& 95\% healthy / 5\% wilted 
& 4839 
& \parbox{7cm}{RGB image data of avocado leaves with texture and color descriptors.\vspace{0.5em}} \\

Seismic Bumps 
& 93\% no bump / 7\% bump 
& 2584 
& \parbox{7cm}{Geophysical and mining operation features from coal mines.\vspace{0.5em}} \\

Bank Marketing 
& 88\% not subscribed / 12\% subscribed 
& 45211 
& \parbox{7cm}{Data from Portuguese bank marketing campaigns, including demographics and contact features.\vspace{0.5em}} \\

Smart Poqueira (own dataset) 
& 99\% visitors / 1\% residents 
& 80,750 
& \parbox{7cm}{Data collected via LPR cameras in a region near Sierra Nevada (Granada, Spain).} \\
\bottomrule
\end{tabular}%
}
\end{table}

\subsection{Experimental setup}

We ran the representative selection methods on an NVIDIA GeForce RTX 3090 GPU with 25.45 GB of memory and CUDA Compute Capability 8.6. Our code executed within an Anaconda environment using Python 3.9.18, and we measured the execution time in each experiment. For training the SSL models, we used the Adam optimizer with a batch size of 32 and a learning rate of $1\times10^{-3}$.

\subsection{Representative Selection Algorithms}

We evaluate the different methods of selecting representative samples used in the literature: Random Sampling and Stratified Sampling \cite{wang2022unsupervised}; centroid-based coverage via K-Means, K-Means$++$, and Bisecting K-Means$++$ that partition the space and pick the nearest point to each centroid \cite{roy2024impact}; USL, which scores samples by inverse $k$-NN distance, clusters into $N$ groups, selects the highest-density exemplar per cluster, and iteratively penalizes duplicates \cite{wang2022unsupervised}; FDMat, which minimizes the maximum mean discrepancy between the sample distributions of the selected subset and the entire dataset \cite{xiao2024feature}; Hybrid-CEAL, which employs entropy-based uncertainty and outlier-based diversity sampling to select high-confidence samples \cite{shankaranarayana2024hybrid}; ESC-FFS, which iteratively selects distant points to broaden subspace coverage and balance class representation in imbalanced datasets \cite{you2020self}; and Gaussian Mapping, an RKA-derived method that aligns samples with minority-class centroids using a Gaussian kernel \cite{li2025kernel}.


\subsection{SSL Algorithms} \label{subsec:algo}

For model training, we use different SSL algorithms commonly applied to tabular datasets in the literature, including: Autoencoder \cite{vincent2008extracting}, VIME \cite{yoon2020vime}, Manifold Mixup \cite{verma2019manifold}, and Contrastive Mixup \cite{darabi2021contrastive}.


\subsection{Evaluation metrics}

We apply each representative selection method using a budget of \(N\) samples. For each selection, we compute the following evaluation metrics of each representative selection method:

\begin{itemize}
  \item \textbf{Imbalance Ratio (IR)} \cite{he2009learning}, to measure class balance among the selected samples:
  \[
    \mathrm{IR} = \frac{\max_{c\in\{0,1\}} N_c}{\min_{c\in\{0,1\}} N_c},
  \]
  where \(N_c\) is the number of selected samples from class \(c\). Values of IR close to one indicate a balanced selection of both classes.

  \item \textbf{Global Fisher Discriminant Ratio (FDR)} \cite{fisher1936use}, to measure class separability in feature space:
  \[
    \mathrm{FDR} = \frac{(\mu_1 - \mu_0)^\top(\mu_1 - \mu_0)}%
    {\displaystyle\sum_{c\in\{0,1\}}\sum_{j=1}^{d}\sigma_{c,j}^2},
  \]
  where \(\mu_c\) is the mean vector and \(\sigma_{c,j}^2\) the variance of feature \(j\) in class \(c\). Higher FDR values reflect greater separation between the two classes.

  \item \textbf{Execution time}, measured in seconds for each selection method.
\end{itemize}

After selecting \(N\) representative samples, we split them into 80–20\% training and validation sets. Each SSL model is trained on the subset of each selection method and evaluated on validation. For evaluating the SSL models trained with each representative selection method we use precision, recall, F1-score, and accuracy~\cite{lever2016classification}. F1-score is used as the primary metric for imbalanced data~\cite{wardhani2019cross}, and weighted averages are reported to account for class frequency imbalance. Finally, we test the model with the original \(M - N\) samples that were not picked by the representation selection method and for which we also have their true labels.

\subsection{Results and discussion}

In this section, we first analyze the internal components and configuration of EDRS through a sensitivity analysis. Subsequently, we benchmark EDRS against state-of-the-art selection methods and, finally we perform an ablation test to measure the contribution of the extremity and density score in EDRS.

\subsubsection{Sensitivity Analysis} \label{subsec:sensitivity}
The performance of EDRS as a preprocessing step relies on four key parameters: the contamination rate for outlier removal ($\%\_{outliers}$), the number of neighbors for density estimation ($k'$), the density quantile threshold ($\delta_{density}$), and the extremity threshold ($e_{thresh}$).

To optimize these parameters effectively, we performed a sensitivity analysis on the representative synthetic dataset (25,400 samples, $\text{FDR}=1$, 90-10 imbalance) with a budget of $N=250$ using Contrastive Mixup. Table~\ref{tab:gridsearch} summarizes the parameters explored during the grid search process and presents the final configuration derived from it. Based on this analysis, the optimal configuration ($\%\_{outliers}=0.10$, $k'=5$, $\delta_{density}=0.10$, $e_{thresh}=0.80$) was applied to all subsequent experiments.

\begin{table}[H]
\centering
\caption{\\ Sensitivity analysis performed on the representative synthetic dataset (25,400 samples, $\text{FDR}=1$, 90-10 imbalance).}
\label{tab:gridsearch}
\resizebox{0.8\textwidth}{!}{%
\begin{tabular}{lcc}
\hline
\textbf{Parameter}                    & \textbf{Search Space}                                                                      & \textbf{Selected}                                      \\ \hline
{ Outliers ($\%\_{outliers}$)}           & { \{0.01, 0.05, 0.10, 0.15, 0.20, 0.25\}}                                       & { \textbf{0.10}}                  \\
{ Neighbors ($k'$)}       & { \{2, 5, 15, 25, 30, 35\}}                              & { \textbf{5}}                   \\
{ Density Quantile ($\delta_{density}$)}          & { \{0.0, 0.10, 0.20, 0.25, 0.30, 0.35\}}                                   & { \textbf{0.10}}                   \\
{ Extremity Thresh. ($e_{thresh}$)}        & { \{0.50, 0.60, 0.70, 0.80, 0.90\}}                               & { \textbf{0.80}}                   \\ \hline
\end{tabular}
}
\end{table}


In Fig.~\ref{fig:threshold_noise}, we analyzed the impact of the extremity threshold ($e_{thresh}$) on model performance (weighted F1-score) and the outlier rate. We define an outlier as any instance located more than 3 standard deviations ($>3\sigma$)  from the majority class centroid. As can be seen, there is a trade-off between minority class coverage and potential outlier propagation.

\begin{figure}[H]
    \centering
    \includegraphics[width=0.85\textwidth]{Figures/threshold_outlier_analysis_final.png} 
    \caption{. Impact of the extremity threshold ($e_{thresh}$) on F1-Score (blue) and Outlier Rate (red) for the representative synthetic dataset. The selected value of 0.80 maximizes performance before outlier propagation degrades the quality of the labeled set.}
    \label{fig:threshold_noise}
\end{figure}

The results indicate an optimal operating point, consistent with the findings obtained in the previous grid search experiment. Lower thresholds ($\le0.70$) fail to trigger the extremity criterion sufficiently, resulting in suboptimal F1-scores. Conversely, overly aggressive thresholds ($\ge0.9$) cause a sharp spike in outlier rates (jumping from 3.9\% to over 9.1\%), which correlates with a decline in classification performance. Hence, the value $e_{thresh}=0.80$ represents the optimal balance, achieving the peak F1-score ($\approx 0.59$) while maintaining a moderate outlier inclusion rate ($\approx 3.9\%$).

To ensure reproducibility across all experiments, we applied a similar optimization methodology to the rest of the methods. The hyperparameters for each comparison method were systematically tuned using a Grid Search strategy. Table~\ref{tab:hyperparams} details the specific search spaces explored and the final parameter values selected for the subsequent methods.

\begin{table}[h!]
\centering
\caption{\\ Hyperparameter search space and final settings for comparison methods.}
\label{tab:hyperparams}
\resizebox{\textwidth}{!}{% Ajusta la tabla al ancho del texto si es necesario
\begin{tabular}{@{}lllll@{}}
\toprule
\textbf{Algorithm} & \textbf{Parameter} & \textbf{Description} & \textbf{Search Space} & \textbf{Selected Value} \\ \midrule
\textbf{K-Means / Variants} & \texttt{init} & Initialization method & $\{\text{random}, \text{k-means++}\}$ & $\text{k-means++}$ \\
\textit{(K-Means, ++, Bisecting)} & \texttt{max\_iter} & Maximum iterations & $\{100, 300, 500\}$ & $300$ \\
 & \texttt{n\_init} & Number of initializations & $\{10, 20\}$ & $10$ \\ \midrule
\textbf{USL} & $k$ & Nearest Neighbors for density & $\{5, 10, 20, 50\}$ & $10$ \\
 & \texttt{metric} & Distance metric & $\{\text{Euclidean}, \text{Cosine}\}$ & $\text{Euclidean}$ \\ \midrule
\textbf{FDMat} & $\lambda$ & Regularization parameter & $\{0.1, 1.0, 10.0\}$ & $1.0$ \\
 & \texttt{transport} & Transport algorithm & $\{\text{Sinkhorn}, \text{EMD}\}$ & $\text{Sinkhorn}$ \\ \midrule
\textbf{Hybrid-CEAL} & \texttt{strategy} & Query strategy & $\{\text{Entropy}, \text{Margin}, \text{LeastConf}\}$ & $\text{Entropy}$ \\
 & \texttt{batch} & Selection batch size & $\{N\}$ (Fixed by budget) & $N$ \\ \midrule
\textbf{ESC-FFS} & $\lambda$ & Sparsity regularization & $\{0.01, 0.1, 1.0\}$ & $0.1$ \\
 & \texttt{max\_iter} & Optimization iterations & $\{100, 500\}$ & $100$ \\ \midrule
\textbf{Gaussian Mapping} & $\sigma$ & Kernel bandwidth & $\{0.1, 0.5, 1.0, 2.0\}$ & $1.0$ \\
 & \texttt{update} & Centroid update rate & $\{0.9, 0.99, 0.999\}$ & $0.99$ \\ \bottomrule
\end{tabular}%
}
\end{table}


\subsubsection{Comparative Analysis}
\label{subsec:results}

We present the results obtained following the pipeline described in Section \ref{sec:method}. First, to evaluate the influence of the datasets' FDR on the performance of the selection methods, Fig.~\ref{fig:synthetic_dataset_1} shows a scatter plot of the distribution of the two classes of each synthetic dataset for the first 2 {\color{blue}principal components}. It shows class separation at the maximum (FDR$=$4) and minimum (FDR $=$ 0.125) thresholds for dataset sizes of 800 (see in Figs.~\ref{fig:menos_4} and \ref{fig:menos_0_125}) and 50,000 (see in Figs.~\ref{fig:mas_4} and \ref{fig:mas_0_125}).

\begin{figure*}[!t]
         \centering
        \begin{subfigure}[b]{0.40\textwidth}
            \centering
            \includegraphics[width=\textwidth]{Figures/menos_4.png}
            \caption[]%
            {{\small  FDR$=$4.0/800 samples}}    
            \label{fig:menos_4}
        \end{subfigure}
        %\hfill 
        \begin{subfigure}[b]{0.40\textwidth}  
            \centering 
            \includegraphics[width=\textwidth]{Figures/menos_0_125.png}
            \caption[]%
            {{\small  FDR$=$0.125/800 samples}}    
            \label{fig:menos_0_125}
        \end{subfigure}
        \begin{subfigure}[b]{0.40\textwidth}
            \centering
            \includegraphics[width=\textwidth]{Figures/mas_4.png}
            \caption[]%
            {{\small  FDR$=$4.0/50,000 samples}}    
            \label{fig:mas_4}
        \end{subfigure}
        %\hfill
    \begin{subfigure}[b]{0.40\textwidth}  
            \centering 
            \includegraphics[width=\textwidth]{Figures/mas_0_125.png}
            \caption[]%
            {{\small  FDR$=$0.125/50,000 samples}}    
            \label{fig:mas_0_125}
        \end{subfigure}
        \caption{. First two PCA for synthetic datasets with 800 (a,b) and 50,000 (c,d) samples, 90-10 imbalance, and FDR$=4,0.125$: class 0 majority (blue); class 1 minority (red).}
        \label{fig:synthetic_dataset_1}
    \end{figure*}

\begin{table}[!t]
\resizebox{\textwidth}{!}{%
\begin{tabular}{cc|cccccccccccc}
\hline
\multicolumn{2}{c|}{FDR}                        & \multicolumn{2}{c|}{4}                    & \multicolumn{2}{c|}{2}                    & \multicolumn{2}{c|}{1}                    & \multicolumn{2}{c|}{0.5}                  & \multicolumn{2}{c|}{0.25}                 & \multicolumn{2}{c}{0.125}        \\ \hline
Methods                                & Budget & IR            & \multicolumn{1}{c|}{Time} & IR            & \multicolumn{1}{c|}{Time} & IR            & \multicolumn{1}{c|}{Time} & IR            & \multicolumn{1}{c|}{Time} & IR            & \multicolumn{1}{c|}{Time} & IR            & Time             \\ \hline
\multirow{3}{*}{Random}                & 100    & 19.00         & \textbf{0.0005}           & 19.00         & \textbf{0.0005}           & 19.00         & \textbf{0.0005}           & 19.00         & \textbf{0.0005}           & 19.00         & \textbf{0.0005}           & 19.00         & \textbf{0.0005} \\
                                       & 250    & 12.16         & \textbf{0.0005}           & 12.16         & \textbf{0.0005}           & 12.16         & \textbf{0.0005}           & 12.16         & \textbf{0.0005}           & 12.16         & \textbf{0.0005}           & 12.16         & \textbf{0.0005} \\
                                       & 500    & 9.42          & \textbf{0.0005}           & 9.42          & \textbf{0.0005}           & 9.42          & \textbf{0.0005}           & 9.42          & \textbf{0.0005}           & 9.42          & \textbf{0.0005}           & 9.42          & \textbf{0.0005} \\ \hline
\multirow{3}{*}{Stratified}            & 100    & 9.00          & 0.0021                    & 9.00          & 0.0027                    & 9.00          & 0.0027                    & 9.00          & 0.0027                    & 9.00          & 0.0027                    & 9.00          & 0.0027          \\
                                       & 250    & 9.00          & 0.0026                    & 9.00          & 0.0026                    & 9.00          & 0.0026                    & 9.00          & 0.0026                    & 9.00          & 0.0026                    & 9.00          & 0.0026          \\
                                       & 500    & 9.00          & 0.0031                    & 9.00          & 0.0031                    & 9.00          & 0.0031                    & 9.00          & 0.0031                    & 9.00          & 0.0031                    & 9.00          & 0.0031          \\ \hline
\multirow{3}{*}{K-Means}               & 100    & 5.25          & 0.3512                    & 4.59          & 0.2193                    & 4.26          & 0.0874                    & 4.11          & 0.0846                    & 4.04          & 0.0833                    & 4.00          & 0.0819          \\
                                       & 250    & 7.93          & 0.1377                    & 7.35          & 0.1287                    & 7.06          & 0.1197                    & 6.79          & 0.1235                    & 6.65          & 0.1280                    & 6.58          & 0.1272          \\
                                       & 500    & 7.93          & 0.2186                    & 8.01          & 0.2073                    & 8.09          & 0.1960                    & 7.93          & 0.2001                    & 7.85          & 0.2022                    & 7.77          & 0.2043          \\ \hline
\multirow{3}{*}{K-Means$++$}           & 100    & 1.13          & 0.0592                    & 1.15          & 0.0593                    & 1.17          & 0.0593                    & 1.30          & 0.0586                    & \textbf{1.29} & 0.0583                    & \textbf{1.33} & 0.0579          \\
                                       & 250    & 2.29          & 0.1304                    & 2.54          & 0.1283                    & 2.79          & 0.1262                    & 2.82          & 0.1268                    & 2.83          & 0.1271                    & 2.85          & 0.1274          \\
                                       & 500    & 5.25          & 0.2417                    & 5.42          & 0.2347                    & 5.58          & 0.2277                    & 5.62          & 0.2319                    & 5.65          & 0.2340                    & 5.67          & 0.2361          \\ \hline
\multirow{3}{*}{Bisecting K-Means$++$} & 100    & 1.27          & 0.3753                    & 1.36          & 0.3679                    & 1.44          & 0.3604                    & 1.47          & 0.3655                    & 1.49          & 0.3670                    & 1.50          & 0.3705          \\
                                       & 250    & 2.52          & 1.3627                    & 2.78          & 1.3778                    & 3.03          & 1.3929                    & 3.00          & 1.3696                    & 2.99          & 1.3580                    & 2.97          & 1.3464          \\
                                       & 500    & 5.25          & 4.2202                    & 5.42          & 4.1360                    & 5.58          & 4.0518                    & 5.54          & 4.1399                    & 5.51          & 4.1839                    & 5.49          & 4.2279          \\ \hline
\multirow{3}{*}{USL}                   & 100    & 1.13          & 0.9240                    & 1.18          & 0.8851                    & 1.22          & 0.8461                    & \textbf{1.25} & 0.7904                    & 1.34          & 0.7625                    & 1.38          & 0.7346          \\
                                       & 250    & 2.12          & 3.4805                    & 2.27          & 3.1497                    & 2.42          & 2.8189                    & \textbf{2.40} & 3.2538                    & \textbf{2.39} & 3.4713                    & \textbf{2.38} & 3.6887          \\
                                       & 500    & 5.25          & 4.1476                    & 5.37          & 4.1501                    & 5.49          & 4.1526                    & \textbf{5.37} & 4.7642                    & \textbf{5.31} & 5.0699                    & \textbf{5.25} & 5.3757          \\ \hline

 \multirow{3}{*}{FDMat}         &  100    &  1.22         &  0.2017                   &  1.25         &  0.2012                 &  1.29          &  0.2015                    &  1.35          &  0.2015                    &  1.40          &  0.2020                    &  1.42          &  0.2021          \\
                                       &  250    &  2.45          &  0.2051                    &  2.58          &  0.2055                    &  2.65          &  0.2061                    &  2.70          &  0.2064                    &  2.75          &  0.2068                    &  2.78          &  0.2072          \\
                                       &  500    &  5.35          &  0.2080                    &  5.48          &  0.2085                    &  5.55          &  0.2091                    &  5.60          &  0.2095                    &  5.63          &  0.2098                    &  5.65          &  0.2102          \\ \hline

 \multirow{3}{*}{Hybrid-CEAL}   &  100    &  1.12          &  0.6520                    &  1.15          &  0.6610                    &  1.18          &  0.6650                    &  1.28          &  0.6680                    &  1.35          &  0.6710                    &  1.41          &  0.6757          \\
                                       &  250    &  2.05          &  0.6853                    &  2.18          &  0.6951                    &  2.30          &  0.7051                    &  2.45          &  0.7153                    &  2.50          &  0.7200                    &  2.55          &  0.7300          \\
                                       &  500    &  5.10          &  0.7020                    &  5.25          &  0.7150                    &  5.40          &  0.7280                    &  5.45          &  0.7403                    &  5.50          &  0.7558                    &  5.58          &  0.7701          \\ \hline

\multirow{3}{*}{ESC-FFS}               & 100    & 0.25          & 2.3645                    & 0.35          & 2.2960                    & 0.45          & 2.2274                    & 0.49          & 2.2510                    & 0.50          & 2.2630                    & 0.52          & 2.2745          \\
                                       & 250    & 2.29          & 19.3666                   & 2.49          & 23.9000                   & 2.68          & 28.4327                   & 2.77          & 23.8850                   & 2.81          & 21.6110                   & 2.85          & 19.3374         \\
                                       & 500    & 5.25          & 80.3449                   & 5.42          & 78.2900                   & 5.58          & 76.2355                   & 5.58          & 79.1990                   & 5.58          & 80.6810                   & 5.58          & 82.1623         \\ \hline

 \multirow{3}{*}{Gaussian Mapping (RKA)} &  100    &  1.10          &  0.3710                    &  1.13          &  0.3715                    &  1.15          &  0.3720                    &  2.80          &  0.3725                    &  3.50          &  0.3730                    &  4.20          &  0.3740          \\
                                       &  250    &  1.95          &  0.3751                    &  2.12          &  0.3763                    &  2.30          &  0.3771                    &  5.10          &  0.3785                    &  6.50          &  0.3803                    &  7.83          &  0.3815          \\
                                       &  500    &  5.05          &  0.3853                    &  5.24          &  0.3882                    &  5.42          &  0.3916                    &  8.50          &  0.3940                    &  11.20         &  0.3974                    &  12.50         &  0.4009          \\ \hline

\multirow{3}{*}{EDRS (ours)}           & 100    & \textbf{1.08} & 0.0666                    & \textbf{1.10} & 0.0509                    & \textbf{1.12} & 0.0351                    & 4.61          & 0.0381                    & 6.35          & 0.0395                    & 8.09          & 0.0410          \\
                                       & 250    & \textbf{1.89} & 0.1171                    & \textbf{2.07} & 0.0907                    & \textbf{2.25} & 0.0642                    & 8.44          & 0.0666                    & 11.53         & 0.0678                    & 14.62         & 0.0690          \\
                                       & 500    & \textbf{4.98} & 0.1451                    & \textbf{5.18} & 0.1292                    & \textbf{5.38} & 0.1132                    & 14.69         & 0.1191                    & 19.35         & 0.1221                    & 24.00         & 0.1250          \\ \hline
\end{tabular}
}
\caption{\\ Summary of sampling performance: for each method and budget, we report the IR and processing time on the synthetic dataset (800 samples, 90–10).}
\label{tab:methods_dataset_synthetic_1}
\end{table}

Table~\ref{tab:methods_dataset_synthetic_1} presents {\color{blue}the} IR and execution time {\color{blue}for small synthetic datasets (800 samples) across} varying FDR values. EDRS achieves the best trade-off when FDR~$\geq$~1, reaching IRs of 1.08–1.12 at budget~100 in under 0.07~s and staying below 0.15~s at budget~500. Random and stratified sampling are fastest (0.0005–0.003~s) but yield poor balance (IRs 9–19). K-Means methods improve IR as FDR decreases, though with higher execution times (0.08–0.35~s at budget~100, increasing with budget). When FDR~$<$~1, USL attains the lowest IR but with higher execution times. ESC-FFS also reaches low IR for small budgets at low FDR, but its execution time grows sharply from about 2.3~s (budget~100) to over 80~s (budget~500). FDMat shows stable execution times (approx. 0.20~s) regardless of the budget, offering moderate IR improvements. Hybrid-CEAL provides competitive IRs comparable to USL but incurs higher computational costs (approx. 0.65–0.77~s). Gaussian mapping demonstrates performance similar to EDRS in environments with high FDR and shows superior robustness (less IR degradation) in scenarios with low FDR ($<$~0.5) compared to EDRS, with constant execution times (approx. 0.37–0.40~s). However, Gaussian mapping does not manage to outperform USL at low FDR ($\leq$~0.5), although it is significantly more competitive in terms of time.

\begin{table}[!t]
\resizebox{\textwidth}{!}{%
\begin{tabular}{cc|cccccccccccc}
\hline
\multicolumn{2}{c|}{FDR}                        & \multicolumn{2}{c|}{4}                    & \multicolumn{2}{c|}{2}                    & \multicolumn{2}{c|}{1}                    & \multicolumn{2}{c|}{0.5}                  & \multicolumn{2}{c|}{0.25}                 & \multicolumn{2}{c}{0.125}        \\ \hline
Methods                                & Budget & IR            & \multicolumn{1}{c|}{Time} & IR            & \multicolumn{1}{c|}{Time} & IR            & \multicolumn{1}{c|}{Time} & IR            & \multicolumn{1}{c|}{Time} & IR            & \multicolumn{1}{c|}{Time} & IR            & Time             \\ \hline
\multirow{3}{*}{Random}                & 100    & 11.50         & \textbf{0.0012}           & 11.50         & \textbf{0.0012}           & 11.50         & \textbf{0.0012}           & 11.50         & \textbf{0.0012}           & 11.50         & \textbf{0.0012}           & 11.50         & \textbf{0.0012} \\
                                       & 250    & 15.67         & \textbf{0.0012}           & 15.67         & \textbf{0.0012}           & 15.67         & \textbf{0.0012}           & 15.67         & \textbf{0.0012}           & 15.67         & \textbf{0.0012}           & 15.67         & \textbf{0.0012} \\
                                       & 500    & 13.29         & \textbf{0.0012}           & 13.29         & \textbf{0.0012}           & 13.29         & \textbf{0.0012}           & 13.29         & \textbf{0.0012}           & 13.29         & \textbf{0.0012}           & 13.29         & \textbf{0.0012} \\ \hline
\multirow{3}{*}{Stratified}            & 100    & 9.00          & 0.0046                    & 9.00          & 0.0046                    & 9.00          & 0.0046                    & 9.00          & 0.0046                    & 9.00          & 0.0046                    & 9.00          & 0.0046          \\
                                       & 250    & 9.00          & 0.0050                    & 9.00          & 0.0050                    & 9.00          & 0.0050                    & 9.00          & 0.0050                    & 9.00          & 0.0050                    & 9.00          & 0.0050          \\
                                       & 500    & 9.00          & 0.0059                    & 9.00          & 0.0059                    & 9.00          & 0.0059                    & 9.00          & 0.0059                    & 9.00          & 0.0059                    & 9.00          & 0.0059          \\ \hline
\multirow{3}{*}{K-Means}               & 100    & 2.23          & 4.0515                    & 1.88          & 3.6767                    & 1.70          & 3.3019                    & 1.55          & 3.1132                    & 1.48          & 3.0188                    & 1.44          & 2.9244          \\
                                       & 250    & 3.63          & 5.7761                    & 2.71          & 6.6022                    & 2.25          & 7.4282                    & 2.20          & 6.5548                    & 2.17          & 6.1180                    & 2.16          & 5.6813          \\
                                       & 500    & 5.76          & 7.7131                    & 4.69          & 9.2461                    & 4.15          & 10.7790                   & 3.85          & 10.1174                   & 3.70          & 9.7866                    & 3.63          & 9.4558          \\ \hline
\multirow{3}{*}{K-Means$++$}           & 100    & 0.98          & 0.3539                    & 0.97          & 0.3523                    & 0.96          & 0.3507                    & 0.96          & 0.3484                    & 0.96          & 0.3472                    & 0.96          & 0.3460          \\
                                       & 250    & 0.92          & 0.8435                    & 0.93          & 0.8188                    & 0.94          & 0.7941                    & 0.96          & 0.7995                    & 0.97          & 0.8021                    & 0.97          & 0.8048          \\
                                       & 500    & 0.90          & 1.6447                    & 0.95          & 1.6760                    & 0.97          & 1.7072                    & \textbf{0.98} & 1.6669                    & \textbf{0.98} & 1.6468                    & \textbf{0.98} & 1.6266          \\ \hline
\multirow{3}{*}{Bisecting K-Means$++$} & 100    & 1.13          & 1.7484                    & 1.16          & 1.7365                    & 1.17          & 1.7246                    & 1.05          & 1.9190                    & 0.99          & 2.0162                    & 0.96          & 2.1134          \\
                                       & 250    & 0.87          & 3.3363                    & 0.87          & 3.3712                    & 0.87          & 3.4060                    & 0.94          & 3.5324                    & 0.98          & 3.5956                    & \textbf{1.00} & 3.6588          \\
                                       & 500    & 0.89          & 7.1743                    & 0.92          & 7.2087                    & 0.94          & 7.2431                    & 0.95          & 7.3517                    & 0.96          & 7.4059                    & 0.96          & 7.4602          \\ \hline
\multirow{3}{*}{USL}                   & 100    & 1.17          & 107.8882                  & 1.17          & 90.5829                   & 1.17          & 73.2775                   & 1.12          & 53.4740                   & 1.09          & 43.5723                   & 1.08          & 33.6705         \\
                                       & 250    & 1.36          & 270.2777                  & 1.13          & 177.4462                  & 1.02          & 84.6146                   & 1.09          & 146.5832                  & 1.12          & 177.5674                  & 1.14          & 208.5517        \\
                                       & 500    & 1.39          & 337.8903                  & 1.13          & 294.4852                  & \textbf{1.00} & 251.0800                  & \textbf{0.98} & 339.8795                  & 0.97          & 384.2793                  & 0.97          & 428.6790        \\ \hline

 \multirow{3}{*}{FDMat}         &  100    &  1.45          &  275.4050                  &  1.58          &  280.1500                  &  1.62          &  282.8060                  &  2.20          &  285.5870                  &  2.55          &  288.1000                  &  2.80          &  290.2500         \\
                                       &  250    &  1.85          &  278.2000                  &  1.95          &  282.4507                  &  2.10          &  285.1000                  &  2.75          &  288.3000                  &  3.10          &  291.5096                  &  3.45          &  295.4280         \\
                                       &  500    &  2.15          &  280.8057                  &  2.30          &  285.2001                  &  2.45          &  289.5017                  &  3.20          &  292.4000                  &  3.65          &  296.6570                  &  3.90          &  301.8028         \\ \hline

 \multirow{3}{*}{Hybrid-CEAL}   &  100    &  1.25          &  210.5010                  &  1.30          &  215.4052                  &  1.45          &  218.6020                  &  1.85          &  222.1050                  &  2.10          &  225.4058                  &  2.45          &  228.8040         \\
                                       &  250    &  1.15          &  212.2053                  &  1.25          &  218.8020                  &  1.35          &  221.9043                  &  1.95          &  225.5021                  &  2.25          &  229.6057                  &  2.53          &  232.1050         \\
                                       &  500    &  1.10          &  215.4020                  &  1.18          &  220.5055                  &  1.25          &  224.1051                  &  1.80          &  228.2025                  &  2.15          &  232.3056                  &  2.40          &  235.4020         \\ \hline

\multirow{3}{*}{ESC-FFS}               & 100    & \textbf{1.00} & 873.7392                  & \textbf{1.00} & 843.8387                  & \textbf{1.00} & 813.9382                  & \textbf{1.00} & 820.4311                  & \textbf{1.00} & 823.6775                  & \textbf{1.00} & 826.9239        \\
                                       & 250    & \textbf{0.98} & 3456.8492                 & \textbf{0.99} & 3424.8493                 & \textbf{0.99} & 3392.8493                 & \textbf{1.00} & 3487.7988                 & \textbf{1.00} & 3535.2736                 & \textbf{1.00} & 3582.7483       \\
                                       & 500    & \textbf{0.92} & 12467.9132                & \textbf{0.96} & 12419.4812                & 0.98          & 12371.0492                & \textbf{0.98} & 12336.1120                & \textbf{0.98} & 12318.6434                & \textbf{0.98} & 12301.1748      \\ \hline

 \multirow{3}{*}{Gaussian Mapping (RKA)} &  100    &  1.05          &  305.5000                  &  1.08          &  310.6500                  &  1.15          &  315.8010                  &  1.55          &  319.9500                  &  1.60          &  322.0000                  &  1.92          &  325.5000         \\
                                       &  250    &  1.02          &  308.5860                  &  1.05          &  312.8000                  &  1.10          &  316.1020                  &  1.35          &  321.4030                  &  1.52          &  326.7000                  &  1.87          &  330.0000         \\
                                       &  500    &  0.98          &  310.5000                  &  1.02          &  315.1030                  &  1.08          &  318.7050                  &  1.22          &  324.3000                  &  1.47          &  328.6100                  &  1.73          &  335.9400         \\ \hline

\multirow{3}{*}{EDRS (ours)}           & 100    & \textbf{1.00}          & 0.7208                    & 0.97          & 0.7287                    & 0.95          & 0.7365                    & 1.57          & 0.7258                    & 1.88          & 0.7205                    & 2.03          & 0.7151          \\
                                       & 250    & 0.94          & 1.1047                    & 0.94          & 1.0878                    & 0.94          & 1.0708                    & 1.67          & 1.0945                    & 2.03          & 1.1063                    & 2.21          & 1.1181          \\
                                       & 500    & 0.90          & 1.8990                    & 0.95          & 1.8774                    & 0.97          & 1.8558                    & 1.59          & 1.8849                    & 1.90          & 1.8994                    & 2.05          & 1.9139          \\ \hline
\end{tabular}
}
\caption{\\ Summary of sampling performance: for each method and budget, we report the IR and processing time on the synthetic dataset (50,000 samples, 90–10).}
\label{tab:methods_dataset_synthetic_2}
\end{table}

Table~\ref{tab:methods_dataset_synthetic_2} summarizes, for large datasets, 50,000 samples,  IR and execution time for each method and budget across varying FDR values. ESC-FFS attains the lowest IR (perfect IR$=1.00$ for budgets 100–500) but with very high execution times (800-12000 seconds). Among other methods, EDRS ranks second when FDR~$\geq$~1, achieving IRs of 1.00–0.95 at budget~100 in under 0.74~s and remaining below 1.90~s at budget~500. Random and stratified sampling are fastest (0.0012–0.0059~s) but poorly balanced (IRs 9.00–15.67). K-Means variants improve IR as FDR decreases, with moderate execution times (3–12~s). When FDR~$<$~1, USL attains good IR results (1.12–1.08 at budget~100) in almost 98\% less time than ESC. FDMat incurs significant execution times (approx. 275–302~s) and struggles to achieve competitive IRs (1.45–3.90). Hybrid-CEAL shows good IR performance (1.10–2.53) but suffers from high computational costs (approx. 210–235~s) due to the iterative nature of the outlier detection on 50k samples. Gaussian mapping offers solid performance, especially in configurations with lower FDRs, where it outperforms EDRS, maintaining IR between 1.22--1.92 in the most difficult scenarios (FDR $\le$ 0.5). However, this robustness comes with a substantially higher computational overhead (approx. 305–336~s) compared to EDRS ($<$~2~s). Overall, for FDR~$\geq$~1, EDRS offers the best IR–time trade-off, as ESC-FFS’s superior IR comes at the highest computational cost.

Once we know the influence of FDR, we calculate the performance of each selection method in the real-world datasets. Table~\ref{tab:methods_dataset} summarizes, for each real dataset, the size of the dataset, original balance classes, FDR, and for each representative selection method the IR and execution time in seconds. Fig.~\ref{fig:real_dataset} shows a scatter plot of the distribution of the two classes of each real dataset for the first 2 PCAs. FDR indicates the separation between classes: lower values imply higher overlap, as in the Seismic Bumps (FDR$=0.10$) and Bank Marketing (FDR$=0.02$) datasets (see in Figs.~\ref{fig:seismic} and~\ref{fig:bank}). We apply each representative selection method to budgets $N$ ranging from 100 to 500 annotated samples. This range, consistent with prior works \cite{wang2022unsupervised,roy2024impact}, represents an annotation cost suitable for evaluating how a small number of labeled samples affects SSL. We limit the range to 100–500 because, for small datasets, selecting more samples nearly covers the entire set, which is meaningless. For large datasets, both IR and processing time change little with higher budgets, so this range provides results that are representative enough for analysis. 

\begin{figure*}
         \centering
        \begin{subfigure}[b]{0.48\textwidth}
            \centering
            \includegraphics[width=\textwidth]{Figures/smart_villages.png}
            \caption[]%
            {{\small  Smart Poqueira Dataset.}}    
            \label{fig:smart_village}
        \end{subfigure}
        \hfill
        \begin{subfigure}[b]{0.48\textwidth}   
            \centering 
            \includegraphics[width=\textwidth]{Figures/cancer_breast.png}
            \caption[]%
            {{\small Breast Cancer Dataset.}}    
            \label{fig:cancer}
        \end{subfigure}
         \vskip\baselineskip   
        \begin{subfigure}[b]{0.48\textwidth}  
            \centering 
            \includegraphics[width=\textwidth]{Figures/thyroid.png}
            \caption[]%
            {{\small  Thyroid Disease Dataset.}}    
            \label{fig:thyroid}
        \end{subfigure}
    \hfill
 \begin{subfigure}[b]{0.48\textwidth}   
            \centering 
            \includegraphics[width=\textwidth]{Figures/wilt.png}
            \caption[]%
            {{\small Wilt Dataset.}}    
            \label{fig:wilt}
        \end{subfigure}
        \vskip\baselineskip   
        \begin{subfigure}[b]{0.48\textwidth}  
            \centering 
            \includegraphics[width=\textwidth]{Figures/seismic.png}
            \caption[]%
            {{\small  Seismic Bumps Dataset.}}    
            \label{fig:seismic}
        \end{subfigure}
    \hfill
 \begin{subfigure}[b]{0.48\textwidth}   
            \centering 
            \includegraphics[width=\textwidth]{Figures/bank.png}
            \caption[]%
            {{\small Bank Marketing Dataset.}}    
            \label{fig:bank}
        \end{subfigure}
        \caption{. First two PCA for the six real datasets: class 0 (blue); class 1 (red).}
        \label{fig:real_dataset}
    \end{figure*}

\begin{table}[!t]
\resizebox{\textwidth}{!}{%
\begin{tabular}{cc|cc|cc|cc|cc|cc|cc}
\hline
\multicolumn{2}{c|}{Dataset Name}               & \multicolumn{2}{c|}{Smart Poqueira} & \multicolumn{2}{c|}{Breast Cancer} & \multicolumn{2}{c|}{Thyroid Disease} & \multicolumn{2}{c|}{Wilt}         & \multicolumn{2}{c|}{Seismic Bumps} & \multicolumn{2}{c}{Bank Marketing} \\
\multicolumn{2}{c|}{Classes Balance}            & \multicolumn{2}{c|}{99-1}           & \multicolumn{2}{c|}{63-37}         & \multicolumn{2}{c|}{92-8}            & \multicolumn{2}{c|}{95-5}         & \multicolumn{2}{c|}{93-7}          & \multicolumn{2}{c}{88-12}          \\
\multicolumn{2}{c|}{FDR}                        & \multicolumn{2}{c|}{3.59}           & \multicolumn{2}{c|}{1.89}          & \multicolumn{2}{c|}{0.45}            & \multicolumn{2}{c|}{0.19}         & \multicolumn{2}{c|}{0.10}          & \multicolumn{2}{c}{0.02}           \\
\multicolumn{2}{c|}{Size}                       & \multicolumn{2}{c|}{80,750}         & \multicolumn{2}{c|}{569}           & \multicolumn{2}{c|}{3,772}           & \multicolumn{2}{c|}{4,839}        & \multicolumn{2}{c|}{2,584}         & \multicolumn{2}{c}{45,211}         \\ \hline
Methods                                & Budget & IR              & Time              & IR              & Time             & IR               & Time             & IR              & Time             & IR              & Time             & IR              & Time             \\ \hline
\multirow{3}{*}{Random}                & 100    & 49.00           & \textbf{0.0019}   & 2.45            & \textbf{0.0013}  & 24.00            & \textbf{0.0019}  & \textbf{15.67} & \textbf{0.0221}  & 19.00           & \textbf{0.0017}  & 9.00            & \textbf{0.0144}  \\
                                       & 250    & 49.00           & \textbf{0.0021}   & 1.94            & \textbf{0.0016}  & 10.36            & \textbf{0.0021}  & 34.71          & \textbf{0.0314}  & 14.63           & \textbf{0.0021}  & 8.61            & \textbf{0.0221}  \\
                                       & 500    & 99.00           & \textbf{0.0029}   & 1.70            & \textbf{0.0015}  & 11.50            & \textbf{0.0026}  & 26.78          & \textbf{0.0346}  & 16.24           & \textbf{0.0082}  & 9.42            & \textbf{0.0280}  \\ \hline
\multirow{3}{*}{Stratified}            & 100    & 99.00           & 0.0090            & 1.70            & 0.0059           & 11.50            & 0.0083           & 19.00          & 0.0717           & 13.29           & 0.0220           & 7.33            & 0.1346           \\
                                       & 250    & 124.00          & 0.0092            & 1.69            & 0.0064           & 12.16            & 0.0256           & \textbf{18.23} & 0.0879           & 14.63           & 0.0359           & 7.62            & 0.2264           \\
                                       & 500    & 99.00           & 0.0096            & 1.69            & 0.0070           & 11.82            & 0.0390           & 17.52          & 0.0951           & 14.15           & 0.0621           & 7.62            & 0.2731           \\ \hline
\multirow{3}{*}{K-Means}               & 100    & 17.25           & 5.6928            & 1.63            & 0.4567           & 7.33             & 31.9409          & 24.00          & 11.4584          & 10.11           & 19.8497          & 4.00            & 241.9715         \\
                                       & 250    & 49.00           & 6.9879            & 1.84            & 0.7567           & 6.14             & 43.6995          & 24.00          & 18.1003          & 12.16           & 20.5799          & 3.81            & 179.3882         \\
                                       & 500    & 61.50           & 7.4519            & 1.63            & 1.6992           & 12.16            & 48.9252          & 20.74          & 22.4629          & 11.20           & 28.1210          & 5.41            & 200.4378         \\ \hline
\multirow{3}{*}{K-Means$++$}           & 100    & 4.26            & 10.1263           & 0.47            & 0.0916           & 2.85             & 2.2952           & 24.00          & 1.5561           & 6.14            & 0.8726           & 3.17   & 3.3108           \\
                                       & 250    & 5.76            & 15.0373           & 0.56            & 0.1704           & 2.62             & 2.5936           & 26.78          & 1.7022           & 5.76            & 1.3085           & 3.03   & 6.2857           \\
                                       & 500    & 4.74            & 31.7062           & 1.37            & 0.2749           & 3.95             & 2.7320           & 21.73          & 1.7805           & 6.25            & 1.6275           & 2.91   & 9.7771           \\ \hline
\multirow{3}{*}{Bisecting K-Means$++$} & 100    & 3.76            & 12.1164           & 0.47            & 0.6362           & 3.00             & 14.2852          & 19.00          & 11.1139          & 6.69            & 13.9085          & 3.35            & 42.4041          \\
                                       & 250    & 4.68            & 23.9680           & 0.67            & 1.3254           & 3.39             & 29.0518          & 30.25          & 18.8844          & 5.76            & 28.9501          & 3.55            & 75.5196          \\
                                       & 500    & 5.02            & 36.4930           & 1.38            & 4.2709           & 4.32             & 56.4298          & 19.00          & 22.2340          & \textbf{5.41}   & 57.1954          & 3.31            & 125.5641         \\ \hline
\multirow{3}{*}{USL}                   & 100    & 4.26            & 86.7506           & 0.45            & 2.6291           & 3.00             & 18.0412          & 49.00          & 4.1280           & \textbf{4.56}   & 8.6000           & 4.00            & 166.9931         \\
                                       & 250    & 5.41            & 239.3468          & 0.56            & 2.7231           & 3.17             & 71.1080          & 26.78          & 9.5308           & \textbf{4.81}   & 20.2641          & 3.72            & 355.7133         \\
                                       & 500    & 4.88            & 483.4320          & 1.38            & 2.8527           & 3.72             & 134.3429         & 25.32          & 122.1568         & 6.46            & 106.0081         & 3.50            & 1671.8234        \\ \hline

 \multirow{3}{*}{FDMat}         &  100    &  15.80           &  855.4020           &  1.95             &  0.2852            &  9.50              &  1.5410            &  28.50            &  2.6850            &  10.50            &  1.0590            &  7.50             &  230.1250          \\
                                       &  250    &  18.20           &  860.8200           &  1.88             &  0.2920            &  8.80              &  1.5750            &  35.20            &  2.7117            &  11.25            &  1.0726            &  7.10             &  235.4530          \\
                                       &  500    &  25.50           &  866.2500           &  1.82             &  0.2952            &  9.20              &  1.6151            &  32.80            &  2.7551            &  12.10            &  1.0909            &  7.80             &  239.5800          \\ \hline

\multirow{3}{*}{Hybrid-CEAL}    &  100    &  3.50            &  512.5010           &  0.85             &  0.4052            &  2.95              &  2.2020            &  26.51            &  4.1050            &  6.20             &  1.6058            &  5.80             &  165.8040          \\
                                       &  250    &  3.85            &  525.2053           &  0.92             &  0.4120            &  3.25              &  2.2543            &  23.12            &  4.1521            &  6.50             &  1.6357            &  6.10             &  172.1050          \\
                                       &  500    &  4.10            &  538.4020           &  1.45             &  0.4255            &  3.80              &  2.3051            &  22.50            &  4.2025            &  6.93             &  1.6856            &  6.50             &  178.4020          \\ \hline

\multirow{3}{*}{ESC-FFS}               & 100    & \textbf{1.22}   & 732.4960          & 0.47            & 7.0903           & 6.69             & 312.4222         & 99.00          & 183.4067         & 6.14            & 300.4281         & \textbf{1.32}   & 1562.2961        \\
                                       & 250    & \textbf{1.40}   & 3349.1886         & 1.55            & 123.8580         & 7.06             & 1125.2226        & 82.33          & 1983.3772        & 5.25            & 1817.6324        & \textbf{1.14}   & 3515.6610        \\
                                       & 500    & \textbf{1.50}   & 12571.8203        & 2.09            & 2056.5102        & 6.94             & 7404.1908        & 26.78          & 7267.8384        & 5.49            & 6403.3125        & \textbf{0.87}   & 12808.1757       \\ \hline

\multirow{3}{*}{Gaussian Mapping (RKA)} &  100    &  2.95            &  920.5023           &  0.65             &  0.3550            &  2.45              &  1.8051            &  23.95            &  3.2050            &  4.97             &  1.2050            &  4.61             &  265.5050          \\
                                        &  250    &  3.15            &  935.4052           &  0.88             &  0.3656            &  2.95              &  1.8552            &  21.58            &  3.3023            &  4.12             &  1.2557            &  4.35             &  271.8055          \\
                                        &  500    &  3.40            &  952.8020           &  1.40             &  0.3751            &  3.45              &  1.9058            &  \textbf{17.42}   &  3.4058            &  5.88             &  1.3059            &  5.72             &  280.4050          \\ \hline

\multirow{3}{*}{EDRS (ours)}           & 100    & 2.44            & 8.8031            & \textbf{0.56}   & 0.0851           & \textbf{2.23}    & 3.0071           & 24.00          & 1.7292           & 5.25            & 0.6940           & 4.88            & 9.5106           \\
                                       & 250    & 2.52            & 17.4461           & \textbf{0.81}   & 0.0939           & \textbf{2.85}    & 3.2426           & 21.73          & 1.7544           & 5.41            & 0.7170           & 5.09            & 11.8332          \\
                                       & 500    & 2.62            & 31.3021           & \textbf{1.32}   & 0.1652           & \textbf{3.27}    & 3.2685           & 17.52 & 1.8594           & 6.14            & 0.8528           & 5.94            & 12.7438          \\ \hline
\end{tabular}
}
\caption{\\ Summary of sampling performance: for each method and budget, we report for each real dataset the IR, and execution time.}
\label{tab:methods_dataset}
\end{table}

The results obtained in Table~\ref{tab:methods_dataset} show that random selection is the fastest but least balanced. Method efficiency depends heavily on FDR and dataset size. On large datasets (Smart Poqueira and Bank Marketing), ESC-FFS consistently achieves IR values closest to 1. FDMat and Hybrid-CEAL incur high computational costs on large datasets (Smart Poqueira and Bank Marketing) without offering superior IRs. Our method, EDRS, attains the best IR on Breast Cancer, Thyroid Disease, and Wilt (for budget $=$ 500), and ranks second on Smart Poqueira while reducing execution time by about 99.3\% compared to ESC-FFS. However, when FDR $<$ 0.19, EDRS is surpassed by Bisecting K-Means$++$, Gaussian Mapping, ESC-FFS, and USL, although our method achieves significantly lower times (reducing execution time by approx. 90-99\% on the Bank Marketing dataset). EDRS performs best on datasets with high FDR values, as the extremity score yields more balanced representative samples.

\begin{table}[!t]
\resizebox{\textwidth}{!}{%
\begin{tabular}{cc|cc|cc|cc|cc|cc|cc}
\hline
\multicolumn{2}{c|}{Dataset Name}    & \multicolumn{12}{c}{Synthetic Distribution}                                                                                                                                                                                                                                                                                                                                                                                                                                                                                                                                                            \\
\multicolumn{2}{c|}{Balance Classes} & \multicolumn{12}{c}{90-10}                                                                                                                                                                                                                                                                                                                                                                                                                                                                                                                                                                             \\
\multicolumn{2}{c|}{FDR}             & \multicolumn{2}{c|}{4}                                 & \multicolumn{2}{c|}{2}                                 & \multicolumn{2}{c|}{1}                                 & \multicolumn{2}{c|}{0.5}                               & \multicolumn{2}{c|}{0.25}                              & \multicolumn{2}{c}{0.125}         \\
\multicolumn{2}{c|}{Size}            & 800             & \multicolumn{1}{l|}{50,000}          & 800             & \multicolumn{1}{l|}{50,000}          & 800             & \multicolumn{1}{l|}{50,000}          & 800             & \multicolumn{1}{l|}{50,000}          & 800             & \multicolumn{1}{l|}{50,000}          & 800             & 50,000          \\ \hline
Methods                   & Budget   & \multicolumn{12}{c}{weighted average F1-score}                                                                                                                                                                                                                                                                                                                                                                                                                                                                                                                                                         \\ \hline
Random                    & 100      & 0.5087          & \multicolumn{1}{l|}{\textbf{0.5934}} & 0.4789          & \multicolumn{1}{l|}{0.5538}          & 0.3581          & \multicolumn{1}{l|}{0.4255}          & 0.2854          & \multicolumn{1}{l|}{0.2667}          & 0.2265          & \multicolumn{1}{l|}{0.2932}          & 0.1874          & 0.2489          \\
                          & 250      & 0.5415          & \multicolumn{1}{l|}{\textbf{0.6260}} & 0.5123          & \multicolumn{1}{l|}{0.5760}          & 0.4181          & \multicolumn{1}{l|}{0.4655}          & 0.3400          & \multicolumn{1}{l|}{0.3100}          & 0.2650          & \multicolumn{1}{l|}{0.3234}          & 0.2100          & 0.2905          \\
                          & 500      & 0.5834          & \multicolumn{1}{l|}{\textbf{0.6627}} & 0.5488          & \multicolumn{1}{l|}{0.6099}          & 0.4611          & \multicolumn{1}{l|}{0.4976}          & 0.3822          & \multicolumn{1}{l|}{0.4098}          & 0.3099          & \multicolumn{1}{l|}{0.3541}          & 0.2427          & 0.3155          \\ \hline
Stratified                & 100      & 0.5499          & \multicolumn{1}{l|}{\textbf{0.5934}} & 0.5099          & \multicolumn{1}{l|}{\textbf{0.5698}} & 0.2899          & \multicolumn{1}{l|}{\textbf{0.4477}} & 0.2476          & \multicolumn{1}{l|}{0.2877}          & 0.1833          & \multicolumn{1}{l|}{0.3076}          & 0.1677          & 0.2544          \\
                          & 250      & 0.5786          & \multicolumn{1}{l|}{0.6146}          & 0.5379          & \multicolumn{1}{l|}{\textbf{0.5921}} & 0.3164          & \multicolumn{1}{l|}{\textbf{0.4863}} & 0.2800          & \multicolumn{1}{l|}{0.3300}          & 0.2074          & \multicolumn{1}{l|}{0.3202}          & 0.1908          & 0.3020          \\
                          & 500      & 0.6033          & \multicolumn{1}{l|}{0.6487}          & 0.5678          & \multicolumn{1}{l|}{\textbf{0.6199}} & 0.3689          & \multicolumn{1}{l|}{\textbf{0.5188}} & 0.3088          & \multicolumn{1}{l|}{0.4088}          & 0.2377          & \multicolumn{1}{l|}{0.3688}          & 0.2054          & 0.3255          \\ \hline
K-Means                   & 100      & 0.4887          & \multicolumn{1}{l|}{0.3565}          & 0.4032          & \multicolumn{1}{l|}{0.2600}          & 0.2588          & \multicolumn{1}{l|}{0.1899}          & 0.2244          & \multicolumn{1}{l|}{0.1699}          & 0.2299          & \multicolumn{1}{l|}{0.1776}          & 0.1988          & 0.1577          \\
                          & 250      & 0.5956          & \multicolumn{1}{l|}{0.3986}          & 0.5143          & \multicolumn{1}{l|}{0.3000}          & 0.2821          & \multicolumn{1}{l|}{0.2118}          & 0.2500          & \multicolumn{1}{l|}{0.2000}          & 0.2822          & \multicolumn{1}{l|}{0.2238}          & 0.2500          & 0.1900          \\
                          & 500      & 0.6599          & \multicolumn{1}{l|}{0.4277}          & 0.5654          & \multicolumn{1}{l|}{0.3500}          & 0.3298          & \multicolumn{1}{l|}{0.2499}          & 0.2898          & \multicolumn{1}{l|}{0.2299}          & 0.2976          & \multicolumn{1}{l|}{0.2399}          & 0.2699          & 0.2098          \\ \hline
K-Means$++$               & 100      & 0.5099          & \multicolumn{1}{l|}{0.0788}          & 0.4520          & \multicolumn{1}{l|}{0.2000}          & 0.3899          & \multicolumn{1}{l|}{0.3099}          & 0.3499          & \multicolumn{1}{l|}{0.2899}          & 0.3298          & \multicolumn{1}{l|}{0.2977}          & 0.2876          & 0.2377          \\
                          & 250      & 0.5631          & \multicolumn{1}{l|}{0.0778}          & 0.5223          & \multicolumn{1}{l|}{0.2900}          & 0.4246          & \multicolumn{1}{l|}{0.3354}          & 0.4000          & \multicolumn{1}{l|}{0.3100}          & 0.3834          & \multicolumn{1}{l|}{0.3000}          & 0.3300          & 0.2500          \\
                          & 500      & 0.8099          & \multicolumn{1}{l|}{0.1288}          & 0.6200          & \multicolumn{1}{l|}{0.3300}          & 0.5699          & \multicolumn{1}{l|}{0.3699}          & 0.4899          & \multicolumn{1}{l|}{0.3199}          & 0.4377          & \multicolumn{1}{l|}{0.3076}          & 0.3877          & 0.2677          \\ \hline
Bisecting K-Means$++$     & 100      & 0.5099          & \multicolumn{1}{l|}{0.1388}          & 0.4520          & \multicolumn{1}{l|}{0.1200}          & 0.1799          & \multicolumn{1}{l|}{0.1299}          & 0.1477          & \multicolumn{1}{l|}{0.1199}          & 0.1099          & \multicolumn{1}{l|}{0.0976}          & 0.0877          & 0.0788          \\
                          & 250      & 0.5592          & \multicolumn{1}{l|}{0.1454}          & 0.5200          & \multicolumn{1}{l|}{0.1400}          & 0.1900          & \multicolumn{1}{l|}{0.1323}          & 0.1500          & \multicolumn{1}{l|}{0.1200}          & 0.1073          & \multicolumn{1}{l|}{0.1087}          & 0.0900          & 0.0811          \\
                          & 500      & 0.5799          & \multicolumn{1}{l|}{0.1699}          & 0.5311          & \multicolumn{1}{l|}{0.1500}          & 0.2299          & \multicolumn{1}{l|}{0.1899}          & 0.1899          & \multicolumn{1}{l|}{0.1699}          & 0.1299          & \multicolumn{1}{l|}{0.1299}          & 0.1077          & 0.1099          \\ \hline
USL                       & 100      & 0.5199          & \multicolumn{1}{l|}{0.1977}          & 0.5600          & \multicolumn{1}{l|}{0.2300}          & 0.4499          & \multicolumn{1}{l|}{0.3699}          & 0.4599          & \multicolumn{1}{l|}{0.2977}          & 0.4099          & \multicolumn{1}{l|}{\textbf{0.3099}} & 0.3299          & \textbf{0.2699} \\
                          & 250      & 0.5528          & \multicolumn{1}{l|}{0.2154}          & 0.5745          & \multicolumn{1}{l|}{0.2500}          & 0.4738          & \multicolumn{1}{l|}{0.4033}          & \textbf{0.5170} & \multicolumn{1}{l|}{0.3540}          & 0.4331          & \multicolumn{1}{l|}{\textbf{0.3400}} & \textbf{0.4028} & \textbf{0.3100} \\
                          & 500      & 0.5999          & \multicolumn{1}{l|}{0.2477}          & 0.6154          & \multicolumn{1}{l|}{0.2700}          & 0.5999          & \multicolumn{1}{l|}{0.4399}          & 0.5499          & \multicolumn{1}{l|}{0.3977}          & \textbf{0.5299} & \multicolumn{1}{l|}{\textbf{0.3699}} & \textbf{0.4677} & \textbf{0.3377} \\ \hline

FDMat           &  100       &  0.5587          &  0.4124                &  0.5153          &  0.4350                &  0.4050          &  0.3851                &  0.3559          &  0.2780                &  0.3150          &  0.2350                &  0.2750          &  0.1980          \\
                          &  250       &  0.6125          &  0.4450                &  0.5480          &  0.4622                &  0.4422          &  0.4155                &  0.3920          &  0.3058                &  0.3483          &  0.2623                &  0.3020          &  0.2254          \\
                          &  500       &  0.6550          &  0.4820                &  0.5858          &  0.4957                &  0.4850          &  0.4380                &  0.4250          &  0.3322                &  0.3853          &  0.2851                &  0.3357          &  0.2520          \\ \hline

 Hybrid-CEAL &  100       &  0.6150          &  0.4951                &  0.5428          &  0.3950                &  0.4850          &  0.3657                &  0.4350          &  0.2850                &  0.3852          &  0.2450                &  0.3150          &  0.2050          \\
                          &  250       &  0.6750          &  0.5350                &  0.5758          &  0.4187                &  0.5350          &  0.3920                &  0.4820          &  0.3350                &  0.4283          &  0.2881                &  0.3552          &  0.2425          \\
                          &  500       &  0.7252          &  0.5720                &  0.6256          &  0.4450                &  0.5720          &  0.4250                &  0.5251          &  0.3822                &  0.4654          &  0.3251                &  0.3952          &  0.2854          \\ \hline
                          
ESC-FFS                   & 100      & 0.4299          & \multicolumn{1}{l|}{0.3877}          & 0.4099          & \multicolumn{1}{l|}{0.3677}          & 0.3099          & \multicolumn{1}{l|}{0.2699}          & 0.2799          & \multicolumn{1}{l|}{0.2475}          & 0.2599          & \multicolumn{1}{l|}{0.2164}          & 0.2077          & 0.1843          \\
                          & 250      & 0.4538          & \multicolumn{1}{l|}{0.4100}          & 0.4300          & \multicolumn{1}{l|}{0.3900}          & 0.3488          & \multicolumn{1}{l|}{0.3000}          & 0.3200          & \multicolumn{1}{l|}{0.2600}          & 0.3032          & \multicolumn{1}{l|}{0.2300}          & 0.2500          & 0.2000          \\
                          & 500      & 0.4899          & \multicolumn{1}{l|}{0.4377}          & 0.4599          & \multicolumn{1}{l|}{0.4177}          & 0.3899          & \multicolumn{1}{l|}{0.3471}          & 0.3499          & \multicolumn{1}{l|}{0.3077}          & 0.3299          & \multicolumn{1}{l|}{0.2327}          & 0.2677          & 0.2264          \\ \hline                        

Gaussian Mapping (RKA)        &  100       &  0.6852          &  0.5452                &  0.5651          &  0.3751                &  0.5017 &  0.3650                &  0.4724          &  0.3253                &  \textbf{0.4250} &  0.2480                &  \textbf{0.3550} &  0.2100          \\
                          &  250       &  0.7425          &  0.5850                &  0.5750          &  0.4050                &  0.5450 &  0.3953                &  0.4950          &  0.3621                &  \textbf{0.4650} &  0.2955                &  0.3751          &  0.2553          \\
                          &  500       &  0.7750          &  0.6150                &  0.6850          &  0.4420                &  0.5932 &  0.4352                &  0.5456          &  0.3952                &  0.5051          &  0.3351                &  0.4352          &  0.3023          \\ \hline

EDRS (ours)               & 100      & \textbf{0.7099} & \multicolumn{1}{l|}{0.5677}          & \textbf{0.5899} & \multicolumn{1}{l|}{0.3877}          & \textbf{0.5199}          & \multicolumn{1}{l|}{0.3702}          & \textbf{0.4799} & \multicolumn{1}{l|}{\textbf{0.3399}} & 0.4171          & \multicolumn{1}{l|}{0.2517}          & 0.3477          & 0.2143          \\
                          & 250      & \textbf{0.7630} & \multicolumn{1}{l|}{0.6016}          & \textbf{0.5900} & \multicolumn{1}{l|}{0.4200}          & \textbf{0.5674}          & \multicolumn{1}{l|}{0.4058}          & 0.5010          & \multicolumn{1}{l|}{\textbf{0.3700}} & 0.4527          & \multicolumn{1}{l|}{0.3100}          & 0.3712          & 0.2600          \\
                          & 500      & \textbf{0.7899} & \multicolumn{1}{l|}{0.6377}          & \textbf{0.7099} & \multicolumn{1}{l|}{0.4577}          & \textbf{0.6099}          & \multicolumn{1}{l|}{0.4471}          & \textbf{0.5599} & \multicolumn{1}{l|}{\textbf{0.4099}} & 0.5090          & \multicolumn{1}{l|}{0.3431}          & 0.4477          & 0.3077          \\ \hline
\end{tabular}
}
\caption{\\ Summary of Contrastive Mixup results for full test set performance: for each method and budget, we report for each synthetic dataset the weighted average F1-score.}
\label{tab:SSL_synt_datasets}
\end{table}

Finally, to evaluate the overall performance of the SSL models depending on the selection algorithm used, we trained SSL models using the samples selected by each method. Table~\ref{tab:SSL_synt_datasets} summarizes the SSL classification results for synthetic datasets. With high FDR values, EDRS achieves the highest F1 across most budgets—especially for the smaller size (800 samples), where it consistently outperforms all methods, even at lower FDRs. In the case of larger datasets (50,000 samples), class separation is amplified: with a high FDR, random and stratified selections can rival or surpass EDRS, as random samples become sufficiently representative. However, as FDR decreases, their performance drops, while EDRS remains stable, only being surpassed by USL and Gaussian Mapping when FDR~$<$~0.5.

Table~\ref{tab:SSL_real_datasets} summarizes the SSL classification results which highlight the influence of FDR on method performance. Among the four algorithms tested, Contrastive Mixup achieved the best results; thus, we report only its performance. For high-FDR datasets (Smart Poqueira, Breast Cancer, and Thyroid Disease), EDRS achieves the highest F1-scores. Although ESC-FFS occasionally performs well (on Breast Cancer), its results are not consistent across all tested datasets. Gaussian Mapping emerges as a robust contender, especially in low-FDR scenarios like Wilt, where it achieves the highest F1-scores (0.8522–0.8756), slightly outperforming EDRS. In other low-FDR datasets (Seismic Bumps), where class separation is weaker, K-Means$++$, Gaussian Mapping and Bisecting K-Means$++$ perform better, though EDRS remains competitive and faster (see Table~\ref{tab:methods_dataset}). For Bank Marketing, the combination of low FDR and large sample size favors Random and Stratified selection, with USL ranking next best.

\begin{table}[!t]
\resizebox{\textwidth}{!}{%
\begin{tabular}{ll|cccccc}
\hline
\multicolumn{2}{c|}{Dataset Name}    & \multicolumn{1}{c|}{Smart Poqueira}  & \multicolumn{1}{c|}{Breast Cancer}   & \multicolumn{1}{c|}{Thyroid Disease} & \multicolumn{1}{c|}{Wilt}            & \multicolumn{1}{c|}{Seismic Bumps}   & Bank Marketing  \\
\multicolumn{2}{c|}{Balance Classes} & \multicolumn{1}{c|}{99-1}            & \multicolumn{1}{c|}{63-37}           & \multicolumn{1}{c|}{92-8}            & \multicolumn{1}{c|}{95-5}            & \multicolumn{1}{c|}{93-7}            & 88-12           \\
\multicolumn{2}{c|}{FDR}             & \multicolumn{1}{c|}{3.59}            & \multicolumn{1}{c|}{1.89}            & \multicolumn{1}{c|}{0.45}            & \multicolumn{1}{c|}{0.19}            & \multicolumn{1}{c|}{0.10}            & 0.02            \\
\multicolumn{2}{c|}{Size}            & \multicolumn{1}{c|}{80,750}          & \multicolumn{1}{c|}{569}             & \multicolumn{1}{c|}{3,772}           & \multicolumn{1}{c|}{4,839}           & \multicolumn{1}{c|}{2,584}           & 45,211          \\ \hline
Methods                   & Budget   & \multicolumn{6}{c}{weighted average F1-score}                                                                                                                                                                                                                                                                                                                                                                                       \\ \hline
Random                    & 100      & \multicolumn{1}{c|}{0.9872}          & \multicolumn{1}{c|}{0.7850}          & \multicolumn{1}{c|}{0.4800}          & \multicolumn{1}{c|}{0.8100}          & \multicolumn{1}{c|}{0.7000}          & 0.7500          \\
                          & 250      & \multicolumn{1}{c|}{0.9871}          & \multicolumn{1}{c|}{0.7992}          & \multicolumn{1}{c|}{0.5038}          & \multicolumn{1}{c|}{0.8278}          & \multicolumn{1}{c|}{0.7223}          & 0.7711          \\
                          & 500      & \multicolumn{1}{c|}{0.9854}          & \multicolumn{1}{c|}{0.8060}          & \multicolumn{1}{c|}{0.5150}          & \multicolumn{1}{c|}{0.8350}          & \multicolumn{1}{c|}{0.7350}          & 0.7800          \\ \hline
Stratified                & 100      & \multicolumn{1}{c|}{0.9857}          & \multicolumn{1}{c|}{0.7900}          & \multicolumn{1}{c|}{0.4848}          & \multicolumn{1}{c|}{0.7100}          & \multicolumn{1}{c|}{0.6800}          & \textbf{0.7700} \\
                          & 250      & \multicolumn{1}{c|}{0.9862}          & \multicolumn{1}{c|}{0.8055}          & \multicolumn{1}{c|}{0.5184}          & \multicolumn{1}{c|}{0.7285}          & \multicolumn{1}{c|}{0.6933}          & \textbf{0.7825} \\
                          & 500      & \multicolumn{1}{c|}{0.9872}          & \multicolumn{1}{c|}{0.8140}          & \multicolumn{1}{c|}{0.5320}          & \multicolumn{1}{c|}{0.7400}          & \multicolumn{1}{c|}{0.7050}          & \textbf{0.7900} \\ \hline
K-Means                   & 100      & \multicolumn{1}{c|}{0.9879}          & \multicolumn{1}{c|}{0.7900}          & \multicolumn{1}{c|}{0.5200}          & \multicolumn{1}{c|}{0.7800}          & \multicolumn{1}{c|}{0.7600}          & 0.5000          \\
                          & 250      & \multicolumn{1}{c|}{0.9881}          & \multicolumn{1}{c|}{0.8055}          & \multicolumn{1}{c|}{0.5356}          & \multicolumn{1}{c|}{0.7960}          & \multicolumn{1}{c|}{0.7725}          & 0.5265          \\
                          & 500      & \multicolumn{1}{c|}{0.9884}          & \multicolumn{1}{c|}{0.8120}          & \multicolumn{1}{c|}{0.5500}          & \multicolumn{1}{c|}{0.8100}          & \multicolumn{1}{c|}{0.7850}          & 0.5400          \\ \hline
K-Means$++$               & 100      & \multicolumn{1}{c|}{0.9886}          & \multicolumn{1}{c|}{0.7950}          & \multicolumn{1}{c|}{0.5000}          & \multicolumn{1}{c|}{0.8200}          & \multicolumn{1}{c|}{0.8350}          & 0.5400          \\
                          & 250      & \multicolumn{1}{c|}{0.9889}          & \multicolumn{1}{c|}{0.8002}          & \multicolumn{1}{c|}{0.5198}          & \multicolumn{1}{c|}{0.8327}          & \multicolumn{1}{c|}{0.8433}          & 0.5570          \\
                          & 500      & \multicolumn{1}{c|}{0.9892}          & \multicolumn{1}{c|}{0.8180}          & \multicolumn{1}{c|}{0.5450}          & \multicolumn{1}{c|}{0.8400}          & \multicolumn{1}{c|}{0.8550}          & 0.5750          \\ \hline
Bisecting K-Means$++$     & 100      & \multicolumn{1}{c|}{0.9889}          & \multicolumn{1}{c|}{0.7980}          & \multicolumn{1}{c|}{0.4800}          & \multicolumn{1}{c|}{0.8500}          & \multicolumn{1}{c|}{\textbf{0.8350}} & 0.3600          \\
                          & 250      & \multicolumn{1}{c|}{0.9891}          & \multicolumn{1}{c|}{0.8108}          & \multicolumn{1}{c|}{0.5000}          & \multicolumn{1}{c|}{0.8668}          & \multicolumn{1}{c|}{\textbf{0.8415}} & 0.3921          \\
                          & 500      & \multicolumn{1}{c|}{0.9893}          & \multicolumn{1}{c|}{0.8150}          & \multicolumn{1}{c|}{0.5100}          & \multicolumn{1}{c|}{0.8750}          & \multicolumn{1}{c|}{\textbf{0.8550}} & 0.4100          \\ \hline
USL                       & 100      & \multicolumn{1}{c|}{0.9890}          & \multicolumn{1}{c|}{0.8030}          & \multicolumn{1}{c|}{0.5204}          & \multicolumn{1}{c|}{0.8500} & \multicolumn{1}{c|}{0.8100}          & 0.5800          \\
                          & 250      & \multicolumn{1}{c|}{0.9892}          & \multicolumn{1}{c|}{0.8126}          & \multicolumn{1}{c|}{0.5345}          & \multicolumn{1}{c|}{0.8620} & \multicolumn{1}{c|}{0.8249}          & 0.6101          \\
                          & 500      & \multicolumn{1}{c|}{0.9894}          & \multicolumn{1}{c|}{0.8200}          & \multicolumn{1}{c|}{\textbf{0.5612}} & \multicolumn{1}{c|}{0.8700} & \multicolumn{1}{c|}{0.8400}          & 0.6350          \\ \hline

FDMat      &  100       & \multicolumn{1}{c|}{0.9880}          & \multicolumn{1}{c|}{0.7980}          & \multicolumn{1}{c|}{0.4950}          & \multicolumn{1}{c|}{0.8150}          & \multicolumn{1}{c|}{0.7850}          &  0.7603          \\
                          &  250       & \multicolumn{1}{c|}{0.9883}          & \multicolumn{1}{c|}{0.8092}          & \multicolumn{1}{c|}{0.5120}          & \multicolumn{1}{c|}{0.8290}          & \multicolumn{1}{c|}{0.7981}          &  0.7753          \\
                          &  500       & \multicolumn{1}{c|}{0.9887}          & \multicolumn{1}{c|}{0.8164}          & \multicolumn{1}{c|}{0.5280}          & \multicolumn{1}{c|}{0.8380}          & \multicolumn{1}{c|}{0.8102}          &  0.7851          \\ \hline

 Hybrid-CEAL &  100       & \multicolumn{1}{c|}{0.9888}          & \multicolumn{1}{c|}{0.8150}          & \multicolumn{1}{c|}{0.5180}          & \multicolumn{1}{c|}{0.8350}          & \multicolumn{1}{c|}{0.8050}          &  0.5500          \\
                          &  250       & \multicolumn{1}{c|}{0.9891}          & \multicolumn{1}{c|}{0.8320}          & \multicolumn{1}{c|}{0.5351}          & \multicolumn{1}{c|}{0.8494}          & \multicolumn{1}{c|}{0.8220}          &  0.5754          \\
                          &  500       & \multicolumn{1}{c|}{0.9895}          & \multicolumn{1}{c|}{0.8480}          & \multicolumn{1}{c|}{0.5522}          & \multicolumn{1}{c|}{0.8615}          & \multicolumn{1}{c|}{0.8353}          &  0.6008          \\ \hline

ESC-FFS                   & 100      & \multicolumn{1}{c|}{0.9890}          & \multicolumn{1}{c|}{\textbf{0.8650}} & \multicolumn{1}{c|}{0.5072}          & \multicolumn{1}{c|}{0.7700}          & \multicolumn{1}{c|}{0.7200}          & 0.3000          \\
                          & 250      & \multicolumn{1}{c|}{0.9889}          & \multicolumn{1}{c|}{\textbf{0.8790}} & \multicolumn{1}{c|}{0.5218}          & \multicolumn{1}{c|}{0.7998}          & \multicolumn{1}{c|}{0.7428}          & 0.3223          \\
                          & 500      & \multicolumn{1}{c|}{0.9888}          & \multicolumn{1}{c|}{\textbf{0.8850}} & \multicolumn{1}{c|}{0.5400}          & \multicolumn{1}{c|}{0.8050}          & \multicolumn{1}{c|}{0.7600}          & 0.3400          \\ \hline

Gaussian Mapping (RKA)        &  100       & \multicolumn{1}{c|}{0.9892}          & \multicolumn{1}{c|}{0.8203}          & \multicolumn{1}{c|}{0.5250}          & \multicolumn{1}{c|}{\textbf{0.8522}} & \multicolumn{1}{c|}{0.8310}          &  0.5602          \\
                          &  250       & \multicolumn{1}{c|}{0.9894}          & \multicolumn{1}{c|}{0.8380}          & \multicolumn{1}{c|}{0.5421}          & \multicolumn{1}{c|}{\textbf{0.8650}} & \multicolumn{1}{c|}{0.8356}          &  0.5883          \\
                          &  500       & \multicolumn{1}{c|}{0.9897}          & \multicolumn{1}{c|}{0.8553}          & \multicolumn{1}{c|}{0.5581}          & \multicolumn{1}{c|}{\textbf{0.8756}} & \multicolumn{1}{c|}{0.8505}          &  0.6006          \\ \hline

EDRS (ours)               & 100      & \multicolumn{1}{c|}{\textbf{0.9894}} & \multicolumn{1}{c|}{0.8250}          & \multicolumn{1}{c|}{\textbf{0.5302}} & \multicolumn{1}{c|}{0.8450}          & \multicolumn{1}{c|}{0.8200}          & 0.5200          \\
                          & 250      & \multicolumn{1}{c|}{\textbf{0.9896}} & \multicolumn{1}{c|}{0.8437}          & \multicolumn{1}{c|}{\textbf{0.5462}} & \multicolumn{1}{c|}{0.8553}          & \multicolumn{1}{c|}{0.8304}          & 0.5339          \\
                          & 500      & \multicolumn{1}{c|}{\textbf{0.9899}} & \multicolumn{1}{c|}{0.8600}          & \multicolumn{1}{c|}{\textbf{0.5612}} & \multicolumn{1}{c|}{0.8700}          & \multicolumn{1}{c|}{0.8450}          & 0.5450          \\ \hline
\end{tabular}
}
\caption{\\ Summary of Contrastive Mixup results for full test set performance: for each method and budget, we report for each real dataset the weighted average F1-score.}
\label{tab:SSL_real_datasets}
\end{table}

\subsubsection{Computational Efficiency Analysis}
\label{subsec:complexity}

To clarify the source of the efficiency improvement in EDRS, we analyze its computational complexity against the state-of-the-art. Let $M$ be the dataset size, $D$ the number of features, and $N$ the selection budget. Existing methods typically suffer from high computational overhead due to global computations. Approaches like USL \cite{wang2022unsupervised}, FDMat \cite{xiao2024feature}, Hybrid-CEAL \cite{shankaranarayana2024hybrid}, and Gaussian Mapping \cite{li2025kernel} require computing full distance, cost, or kernel matrices across the entire dataset, leading to a quadratic complexity of $\mathcal{O}(M^2)$. Even more computationally demanding is ESC-FFS \cite{you2020self}, which employs an iterative forward selection method that requires re-evaluating errors over the entire dataset for each selected sample, scaling polynomially as $\mathcal{O}(N \cdot M^2)$. In contrast, the efficiency of EDRS stems from performing density estimation and clustering in a reduced 2D space. By projecting data via PCA ($\mathcal{O}(M \cdot D^2)$) and leveraging spatial indexing structures (k-d trees) for the $k$-nearest neighbor search, we reduce the complexity of the density estimation step to quasi-linear $\mathcal{O}(M \log M)$. Finally, the clustering step scales linearly as $\mathcal{O}(M \cdot N \cdot i)$, making the method highly scalable.

Table~\ref{tab:scalability_analysis} on the right-hand side presents our empirical scalability test. EDRS exhibits near-linear growth, processing 50,000 samples in just 1.37 seconds. In comparison, quadratic methods like FDMat or USL require between 200 and 400 seconds, while ESC-FFS suffers a polynomial explosion, taking over 10 hours for the same task.

\begin{table}[h]
\centering
\caption{\\ \textbf{Left:} Scalability comparison (Mean execution time in seconds) across dataset sizes ($M$). \textbf{Right:} Detailed runtime breakdown (Mean $\pm$ Std. Dev.) of EDRS internal components for $M=50,000$.}
\label{tab:scalability_analysis}
\resizebox{\textwidth}{!}{%
\begin{tabular}{lrrrrr|lr}
\toprule
\multicolumn{6}{c|}{\textbf{Scalability Comparison (Total Time)}} & \multicolumn{2}{c}{\textbf{EDRS Component Breakdown ($M$=50k)}} \\ \cmidrule(r){1-6} \cmidrule(l){7-8} 
\textbf{Method} & \textbf{1k} & \textbf{5k} & \textbf{10k} & \textbf{25k} & \textbf{50k} & \textbf{Component} & \textbf{Time (s)} \\ \midrule
USL & 1.13 & 4.46 & 14.66 & 85.46 & 336.80 & PCA Reduction & $0.11 \pm 0.01$ \\
FDMat & 3.15 & 6.74 & 18.12 & 97.18 & 376.50 & Score Calc. & $0.10 \pm 0.01$ \\
Hybrid-CEAL & 1.59 & 3.60 & 10.02 & 54.86 & 213.08 & Clustering & $1.16 \pm 0.05$ \\
ESC-FFS & 8.06 & 153.02 & 787.63 & 7,149.63 & 37,850.38 & & \\
Gaussian Mapping & 2.11 & 4.88 & 13.52 & 74.01 & 290.88 & & \\
\textbf{EDRS (Ours)} & \textbf{0.10} & \textbf{0.17} & \textbf{0.28} & \textbf{0.67} & \textbf{1.37} & \textbf{Total} & $\mathbf{1.37 \pm 0.06}$ \\ \bottomrule
\end{tabular}%
}
\end{table}

Decomposing the EDRS runtime ($M=50,000$) reveals that the clustering step consumes $\approx$85\% of the total time, whereas the proposed Density and Extremity scores calculation account for only $\approx$15\% (See left-hand side of Table~\ref{tab:scalability_analysis}). This confirms that our composite metric introduces negligible overhead compared to the necessary clustering step, ensuring scalability for large datasets.

\subsubsection{Ablation Study}
To assess the individual contributions of the proposed scoring mechanisms, we conducted an ablation study comparing:
\begin{enumerate}
    \item \textbf{Density-only}: Selecting samples solely based on the density score ($d_i$).
    \item \textbf{Extremity-only}: Selecting samples solely based on the extremity score ($e_i$).
    \item \textbf{EDRS (Combined)}: The proposed method using the composite score $s_i = d_i \cdot (1 + e_i)$.
\end{enumerate}

%We maintained consistent experimental conditions to ensure robust evaluation. 
We maintained consistent experimental conditions to control for confounding factors. We utilized Contrastive Mixup for SSL training, as it yielded optimal validation performance. The SSL budget was fixed at $N=250$. For different ablation tests, Table \ref{tab:all_datasets_combined} presents the IR and F1-score for the twelve synthetic datasets and the six real datasets.

The results indicate that performance depends on class separability. These results are in line with those reported in Section \ref{subsec:results}. In scenarios with FDR $\ge 1$, the extremity-only method outperforms the density-only method. Furthermore, the combined score improves both F1-score and IR in these configurations. Conversely, in scenarios with FDR $< 1$, the density-only metric typically yields superior IR values. The Thyroid Disease and Wilt datasets constitute exceptions; here, the combined metric (density+extremity) also improves IR slightly. Regarding F1-score, the combined metric provides gains in specific instances where FDR $< 1$. This occurs in 800-sample synthetic datasets with FDR $0.5$ and $0.25$ and in 50,000-sample synthetic datasets with FDR $0.5$ (see in Table \ref{tab:all_datasets_combined}). The combined metric also yields minor F1-score improvements for the Thyroid Disease and Seismic Bumps datasets (see in Table \ref{tab:all_datasets_combined}).

\begin{table}[!h]
\centering
\caption{\\ Comparison of Density-only, Extremity-only, and EDRS (Combined) across synthetic and real-world datasets. Experiments performed with Contrastive Mixup ($N=250$).}
\label{tab:all_datasets_combined}
\resizebox{1.0\textwidth}{!}{%
\begin{tabular}{llcccccc}
\toprule
\multirow{2}{*}{\textbf{Dataset Name}} & \multirow{2}{*}{\textbf{FDR}} & \multicolumn{2}{c}{\textbf{Density-only}} & \multicolumn{2}{c}{\textbf{Extremity-only}} & \multicolumn{2}{c}{\textbf{EDRS (Combined)}} \\ \cmidrule(lr){3-4} \cmidrule(lr){5-6} \cmidrule(lr){7-8} 
 &  & \textbf{IR} & \textbf{F1-Score} & \textbf{IR} & \textbf{F1-Score} & \textbf{IR} & \textbf{F1-Score} \\ \midrule
 
% --- SYNTHETIC N=800 ---
% Usamos un tabular anidado para forzar el salto de línea dentro del multirow
\multirow{6}{*}{\begin{tabular}{@{}l@{}}\textbf{Synthetic}\\\textbf{800 samples}\end{tabular}} 
 & $4$ & $2.05$ & $0.5781$ & $1.86$ & $0.7012$ & $\mathbf{1.89}$ & $\mathbf{0.7630}$ \\
 & $2$ & $2.18$ & $0.5821$ & $2.02$ & $0.5865$ & $\mathbf{2.07}$ & $\mathbf{0.5900}$ \\
 & $1$ & $2.37$ & $0.5074$ & $2.15$ & $0.5400$ & $\mathbf{2.25}$ & $\mathbf{0.5674}$ \\ 
 & $0.5$ & $\mathbf{2.51}$ & $0.4965$ & $4.10$ & $0.4810$ & $8.44$ & $\mathbf{0.5010}$ \\
 & $0.25$ & $\mathbf{2.68}$ & $0.4181$ & $6.20$ & $0.4050$ & $11.53$ & $\mathbf{0.4527}$ \\
 & $0.125$ & $\mathbf{3.01}$ & $\mathbf{0.3873}$ & $9.50$ & $0.3620$ & $14.62$ & $0.3712$ \\ \midrule

% --- SYNTHETIC N=50k ---
\multirow{6}{*}{\begin{tabular}{@{}l@{}}\textbf{Synthetic}\\\textbf{50,000 samples}\end{tabular}} 
 & $4$ & $1.12$ & $0.3105$ & $0.92$ & $0.5210$ & $\mathbf{0.94}$ & $\mathbf{0.6016}$ \\
 & $2$ & $1.08$ & $0.2873$ & $0.91$ & $0.3850$ & $\mathbf{0.94}$ & $\mathbf{0.4200}$ \\
 & $1$ & $1.02$ & $0.4021$ & $0.91$ & $0.4032$ & $\mathbf{0.94}$ & $\mathbf{0.4058}$ \\ 
 & $0.5$ & $\mathbf{1.37}$ & $0.3382$ & $1.52$ & $0.3250$ & $1.67$ & $\mathbf{0.3700}$ \\
 & $0.25$ & $\mathbf{1.34}$ & $\mathbf{0.3217}$ & $1.85$ & $0.2900$ & $2.03$ & $0.3100$ \\
 & $0.125$ & $\mathbf{1.68}$ & $\mathbf{0.2998}$ & $2.05$ & $0.2450$ & $2.21$ & $0.2600$ \\ \midrule

% --- REAL WORLD ---
\textbf{Smart Poqueira} & $3.59$ & $4.32$ & $0.9892$ & $3.10$ & $0.9894$ & $\mathbf{2.52}$ & $\mathbf{0.9896}$ \\
\textbf{Breast Cancer} & $1.89$ & $0.49$ & $0.8213$ & $0.62$ & $0.8350$ & $\mathbf{0.81}$ & $\mathbf{0.8437}$ \\
\textbf{Thyroid Disease} & $0.45$ & $2.98$ & $0.5401$ & $3.05$ & $0.5300$ & $\mathbf{2.85}$ & $\mathbf{0.5462}$ \\
\textbf{Wilt} & $0.19$ & $26.92$ & $\mathbf{0.8591}$ & $24.50$ & $0.8500$ & $\mathbf{21.73}$ & $0.8553$ \\
\textbf{Seismic Bumps} & $0.10$ & $\mathbf{5.12}$ & $0.8297$ & $5.30$ & $0.8050$ & $5.41$ & $\mathbf{0.8304}$ \\
\textbf{Bank Marketing} & $0.02$ & $\mathbf{4.22}$ & $\mathbf{0.6078}$ & $4.65$ & $0.5510$ & $5.09$ & $0.5339$ \\ \bottomrule
\end{tabular}%
}
\end{table}


\subsection{Discussion}

EDRS is positioned as the optimal choice for scenarios requiring real-time efficiency or where there is at least a slight degree of separability ($FDR \ge 0.19$). In these contexts, the method demonstrates high efficiency and effectiveness. Specifically, it achieves superior IR and F1-scores, as observed in the Smart Poqueira and Thyroid Disease datasets or synthetic scenarios with $FDR \ge 1$. This performance provides a significant advantage over iterative methods (such as ESC-FFS, USL) or kernel-based approaches (such as Gaussian Mapping). While those alternatives may effectively handle overlap, {\color{blue}they take} drastically {\color{blue}more time (EDRS is 90-99\% faster)}, as shown in Table~\ref{tab:methods_dataset}. Despite these strengths, our experimental analysis identifies specific boundaries regarding its applicability. The primary limitation arises in scenarios with extremely low class separability ($FDR < 0.19$), as observed in the Seismic Bumps ($FDR=0.10$) and Bank Marketing ($FDR=0.02$) datasets. In these cases, our method struggles to prioritize minority class samples, resulting in a degradation of the IR. This pattern of performance drop is consistent with our synthetic experiments (Tables~\ref{tab:methods_dataset_synthetic_1} and \ref{tab:methods_dataset_synthetic_2}), where EDRS showed a degradation in IR at FDR levels $\le 0.5$, being outperformed by methods such as USL or Gaussian Mapping. In these highly overlapping environments, the fundamental assumption of our extremity scoring is weakened. When classes are highly intertwined, extreme values in the linear projection may not correspond to the minority class, but rather to noisy instances of the majority class. Consequently, computationally heavier methods may be preferable only for these specific, highly overlapping problems. 

Our methodology could be extended to datasets containing non-numeric features. Future work could employ proper encoding strategies, such as One-Hot Encoding, to handle nominal features without imposing arbitrary ordinality. This transformation maps categorical states to orthogonal binary vectors, ensuring compatibility with Euclidean distance metrics \cite{hancock2020survey}.  Specifically, this approach ensures that the Euclidean distance between any two distinct categories remains constant, allowing the density estimation step to treat categorical dissimilarities fairly in an equidistant space. Furthermore, combining this encoding with Z-score normalization for numerical features ensures that all dimensions contribute equally to the distance and variance computations.

\section{Conclusion} \label{sec:conclusion} 

In this paper, we present EDRS, an unsupervised method for selecting representative binary samples in SSL on tabular data. EDRS combines density and extremity scores to capture both dense regions and extreme cases, creating representative balanced labeled sets. To our knowledge, it is the first approach to integrate both criteria in an unsupervised pre-training step, compatible with any classifier or SSL model. We evaluate EDRS on twelve synthetic and six real-world tabular datasets with varying class separability (FDR metric). Compared to ten representative selection methods across four SSL algorithms, EDRS achieves balanced binary-class selections (IR $\approx$ 1) for datasets with FDR $\geq$ 1, performing best on small and medium datasets (500-5,000 samples). While the kernel-based Gaussian Mapping method demonstrates superior robustness in scenarios with high class overlap (FDR $<$ 0.5), EDRS consistently outperforms it when separability is moderate to high, effectively targeting the decision boundaries with significantly lower computational overhead.

For larger datasets (40-80k), ESC-FFS slightly improves balance, but EDRS achieves similar results with 99\% faster selection times. Furthermore, EDRS offers a significantly better trade-off between performance and computational cost compared to iterative active learning methods such as Hybrid-CEAL or density-based methods like FDMat, surpassing both in scalability on large datasets while delivering superior selection quality. When combined with Contrastive Mixup, on synthetic data, EDRS performs best on small datasets and remains competitive on larger ones, offering major efficiency gains. On real-world datasets, EDRS improves average F1-scores by 3-5\% over other methods for FDR $>$ 0.45. For datasets exceeding 40k samples, random or stratified sampling {\color{blue}methods} obtain superior results in around 50\% of cases, which means they can serve as practical alternatives when speed is the main priority.


\section*{CRediT authorship contribution statement}
\textbf{Alberto Durán-López:} Methodology, Validation, Investigation, Writing- Original Draft. \textbf{Daniel Bolaños-Martinez:} Methodology, Validation, Investigation, Writing- Review \& Editing. \textbf{Maria Bermudez-Edo:} Investigation, Resources, Writing- Review \& Editing, Supervision, Project administration.

\section*{Declaration of competing interest}
The authors declare that they have no known competing financial interests or personal relationships that could have appeared to influence the work reported in this paper.

\section*{Acknowledgements}

Grant C-SEJ-128-UGR23 funded by Consejería de Universidad, Investigación e Innovación and by ERDF Andalusia Program 2021–2027; project PID2023-149185OB-I00 funded by MICIU/AEI/10.13039/501100011033 and co-funded by ERDF/EU.

%% The Appendices part is started with the command \appendix;
%% appendix sections are then done as normal sections
%\appendix



%\newpage

\section*{Data Availability}

All data used in this paper can be found in \url{https://zenodo.org/records/14262136}

%% If you have bib database file and want bibtex to generate the
%% bibitems, please use
%%
 \bibliographystyle{elsarticle-num}
 \bibliography{./els-biblio}

%% else use the following coding to input the bibitems directly in the
%% TeX file.

%% Refer following link for more details about bibliography and citations.
%% https://en.wikibooks.org/wiki/LaTeX/Bibliography_Management


\end{document}

\endinput
%%
%% End of file `elsarticle-template-num.tex'.
